\section{Étude des fibres d'un morphisme universellement ouvert}
\label{section:IV.15-fr}

\subsection{Multiplicités des fibres d'un morphisme universellement ouvert}
\label{subsection:IV.15.1-fr}

\begin{proposition}[15.1.1]
\label{IV.15.1.1-fr}
\oldpage[IV-3]{223}
Soient $Y$ un préschéma irréductible localement noethérien, de point générique
$\eta$, $f:X\to Y$ un morphisme localement de type fini, $y$ un point de $Y$,
$\mathcal F$ un $\mathcal O_X$-Module cohérent~; on pose
$\mathcal F_y=\mathcal F\otimes_{\mathcal O_Y}\mathbf{k}(y)$ (faisceau de modules
sur la fibre $f^{-1}(y)$). Soit $z$ le point générique d'une composante
irréductible de $\operatorname{Supp}(\mathcal F_y)$, et soient $X_i$
($1\leqslant i\leqslant n$) ceux des sous-préschémas fermés réduits de $X$ dont
l'espace sous-jacent est une composante irréductible de
$\operatorname{Supp}(\mathcal F)$ contenant $z$, et qui sont tels que la
restriction $X_i\to Y$ de $f$ soit un morphisme universellement ouvert au point
$z$~; soit $z_i$ le point générique de $X_i$ ($1\leqslant i\leqslant n$). Si
$\lambda_x(\mathcal F_y)$ (resp. $\lambda_t(\mathcal F_\eta)$) désigne la longueur
géométrique de $\mathcal F_y$ en un point $x\in f^{-1}(y)$ (resp. celle de
$\mathcal F_\eta$ en un point $t\in f^{-1}(\eta)$) (4.7.5), on a
\[
\tag{15.1.1.1}
\label{IV.15.1.1.1-fr}
  \lambda_z(\mathcal F_y)\geqslant
  \sum_{i=1}^{n}\lambda_{z_i}(\mathcal F_\eta).
\]
Si de plus on suppose l'anneau $\mathcal O_y$ régulier, on a
\[
\tag{15.1.1.2}
\label{IV.15.1.1.2-fr}
  \operatorname{long}((\mathcal F_y)_z)\geqslant
  \sum_{i=1}^{n}\operatorname{long}((\mathcal F_\eta)_{z_i}).
\]
\end{proposition}

\begin{proof}
\oldpage[IV-3]{224}
Les fibres $f^{-1}(y')$ de $f$ (pour tout $y'\in Y$) ne changeant pas lorsque
l'on remplace $f$ par $f_{\mathrm{red}}$, on peut supposer que $X$ et $Y$ sont
réduits (2.4.3, (vi)), donc $Y$ intègre~; on posera
$K=\mathbf{k}(\eta)$. On peut en outre remplacer $f$ par la restriction
$\operatorname{Supp}(\mathcal F)\to Y$ de $f$ au sous-préschéma fermé réduit de
$X$ ayant $\operatorname{Supp}(\mathcal F)$ comme espace sous-jacent, donc se
borner au cas où $X=\operatorname{Supp}(\mathcal F)$. Enfin, si $y=\eta$, il n'y
a qu'une seule composante irréductible $X_i$ de $X$ contenant $z$
($\mathbf{0}_{\mathrm I}$, 2.1.8) et on a $z_i=z$, ce qui rend (15.1.1.1) et
(15.1.1.2) triviales (sans hypothèses sur $f$). On peut donc se borner au cas
$y\ne\eta$~; il résulte alors de l'hypothèse et de (1.10.3) que les $z_i$
appartiennent à $f^{-1}(\eta)$, les seconds membres de (15.1.1.1) et
(15.1.1.2) étant donc définis. Soit $Z_i$ le sous-préschéma fermé réduit de
$f^{-1}(\eta)$ ayant pour espace sous-jacent $X_i\cap f^{-1}(\eta)$~; on sait
(4.6.6) qu'il existe une extension finie radicielle $K'$ de $K$ telle que, pour
tout $i$, le $K'$-préschéma $(Z_i\otimes_K K')_{\mathrm{red}}$ soit
géométriquement réduit sur $K'$. En outre, le morphisme projection
$f^{-1}(\eta)\otimes_K K'\to f^{-1}(\eta)$ est fini, dominant et radiciel
(\textbf{I}, 3.5.7), donc est un homéomorphisme (2.4.5). Si $z_i'$ est l'unique
point de $f^{-1}(\eta)\otimes_K K'$ au-dessus de $z_i$, il résulte de (4.7.9) et
(4.7.5) que l'on a
\[
\tag{15.1.1.3}
\label{IV.15.1.1.3-fr}
  \lambda_{z_i}(\mathcal F_\eta)
  =\lambda_{z_i'}(\mathcal F_\eta\otimes_K K')
  =\operatorname{long}((\mathcal F_\eta\otimes_K K')_{z_i'}).
\]

Soit $V$ un anneau de valuation discrète, de corps de fractions $K'$, dominant
$\mathcal O_y$ (\textbf{II}, 7.1.7)~; posons $Y'=\operatorname{Spec}(V)$,
$X'=X\times_Y Y'$, $\mathcal F'=\mathcal F\otimes_Y Y'$,
$f'=f_{(Y')}:X'\to Y'$~; si $g:Y'\to Y$ est le morphisme structural, $\eta'$ le
point générique de $Y'$ et $y'$ son point fermé, on a $g(y')=y$, $g(\eta')=\eta$~;
le morphisme $\operatorname{Spec}(\mathbf{k}(y'))\to
\operatorname{Spec}(\mathbf{k}(y))$ étant fidèlement plat, il en est de même de
la projection $g':f'^{-1}(y')\to f^{-1}(y)$ (2.2.13), donc (2.3.4), il y a un
point générique $z'$ d'une composante irréductible de $f'^{-1}(y')$ tel que
$g'(z')=z$. Par construction, on a $\mathbf{k}(\eta')=K'$, donc
$f'^{-1}(\eta')=f^{-1}(\eta)\otimes_K K'$~; d'autre part, si l'on pose
$X_i'=X_i\times_Y Y'$, la restriction $X_i'\to Y'$ de $f'$ est un morphisme
universellement ouvert au point $z'$, donc (1.10.3) il existe un point
$z_i''\in f'^{-1}(\eta')$ qui est une générisation de $z'$, et l'on peut
évidemment supposer que $z_i''$ est point générique d'une composante
irréductible de $X_i'\cap f'^{-1}(\eta')$~; comme la projection de
$X_i'\cap f'^{-1}(\eta')$ dans $f^{-1}(\eta)$ est $Z_i$, on a $z_i''=z_i'$.
En vertu de (4.7.9) et (4.7.5), on a
\[
\tag{15.1.1.4}
\label{IV.15.1.1.4-fr}
  \lambda_z(\mathcal F_y)=\lambda_{z'}(\mathcal F'_{y'})
  \geqslant\operatorname{long}((\mathcal F'_{y'})_{z'})
\]
et il résulte donc de cette inégalité et de (15.1.1.3) qu'il suffit, pour établir
(15.1.1.1), de démontrer l'inégalité
\[
\tag{15.1.1.5}
\label{IV.15.1.1.5-fr}
  \operatorname{long}((\mathcal F'_{y'})_{z'})\geqslant
  \sum_{i=1}^{n}\operatorname{long}((\mathcal F'_{\eta'})_{z_i'}).
\]

Considérons maintenant la seconde inégalité (15.1.1.2)~; nous utiliserons le
lemme suivant~:

\begin{lemma}[15.1.1.6]
\label{IV.15.1.1.6-fr}
Soient $A$ un anneau local régulier qui n'est pas un corps, $K$ son corps des
fractions, $k$ son corps résiduel. Il existe un anneau de valuation discrète $V$
dominant $A$, ayant $K$ pour corps des fractions et dont le corps résiduel est
une extension transcendante pure de $k$.
\end{lemma}

\oldpage[IV-3]{225}
Posons $S=\operatorname{Spec}(A)$, et considérons le $S$-schéma $P$ obtenu en
faisant éclater le point fermé $s$ de $S$ (\textbf{II}, 8.1.3)~; si $\mathfrak m$
est l'idéal maximal de $A$, on a donc par définition
$P=\operatorname{Proj}(B)$, où $B$ est la $A$-algèbre graduée
$\bigoplus_{n\geqslant0}\mathfrak m^n$. Si $h:P\to S$ est le morphisme
structural, la fibre $h^{-1}(s)$ est isomorphe à
$\operatorname{Proj}(B\otimes_A k)$ (\textbf{II}, 2.8.10), et par définition
$B\otimes_A k$ est la $k$-algèbre graduée
$\bigoplus_{n\geqslant0}\mathfrak m^n/\mathfrak m^{n+1}$~; mais comme $A$ est
régulier, cette algèbre est isomorphe à une algèbre de polynômes
$k[t_1,\ldots,t_e]$ ($t_i$ indéterminées) ($\mathbf{0}_{\mathrm{IV}}$, 17.1.1)
et par suite $h^{-1}(s)$ est isomorphe à $\mathbf P_k^{e-1}$. Soit $\zeta$ le
point générique de $h^{-1}(s)$~; si $t_i'$ ($1\leqslant i\leqslant e$) sont les
éléments de $\mathfrak m$ dont les classes mod. $\mathfrak m^2$ sont les $t_i$,
$B_{(t_e')}$ est l'anneau d'un voisinage ouvert affine de $\zeta$ dans $P$, et
l'on a $\mathbf{k}(\zeta)=k(t_1,\ldots,t_{e-1})$~; d'autre part, comme $A$ est
intégralement clos, on vérifie aisément qu'il en est de même de $B$ (Bourbaki,
\emph{Alg. comm.}, chap.~V, \S~1, n\textsuperscript{o}~8, cor.~1 de la prop.~21),
donc $B_{(t_e')}$ est aussi intégralement clos, et il en est par suite de même de
l'anneau local $\mathcal O_{P,\zeta}$. Enfin, comme $A=\mathcal O_s$ est
régulier, donc un anneau universellement caténaire (5.6.4), on a, en vertu de
(5.6.5),
$\dim(\mathcal O_{P,\zeta})=\dim(\mathcal O_s)-(e-1)$, puisque $h$ est un
morphisme birationnel, que $\dim(h^{-1}(s))=e-1$ et que l'anneau local de la
fibre $h^{-1}(s)$ en son point générique $\zeta$ est de dimension $0$. Mais
$\dim(\mathcal O_s)=e$, donc $\dim(\mathcal O_{P,\zeta})=1$, et
$\mathcal O_{P,\zeta}=V$ est un anneau de valuation discrète, étant noethérien,
de dimension $1$ et intégralement clos (\textbf{II}, 7.1.6)~; il répond
évidemment à la question.

Ce lemme étant établi, on peut reprendre le raisonnement fait pour l'inégalité
(15.1.1.1) avec $Y'=\operatorname{Spec}(V)$~; cette fois $\mathbf{k}(y')$ est
une extension \emph{séparable} de $\mathbf{k}(y)$, et par suite (4.7.9), on a
$\operatorname{long}((\mathcal F_y)_z)=
\operatorname{long}((\mathcal F'_{y'})_{z'})$~; comme par ailleurs
$f'^{-1}(\eta')=f^{-1}(\eta)$ et $\mathcal F'_{\eta'}=\mathcal F_\eta$, on est
encore ramené à prouver l'inégalité (15.1.1.5). Or, si $t$ est une uniformisante
de $\mathcal O_{y'}$, $f'^{-1}(y')$ est le sous-préschéma fermé de $X'$ défini
par l'idéal $t\mathcal O_{X'}$, et on a
$\mathcal F'_{y'}=\mathcal F'/t\mathcal F'$~; d'autre part, on vérifie aussitôt
(\textbf{I}, 9.1.12) que
$(\mathcal F'_{\eta'})_{z_i'}=\mathcal F'_{z_i'}$~; l'inégalité à démontrer
(15.1.1.5) n'est autre alors qu'un cas particulier de (3.4.1.1).
\end{proof}

\begin{corollary}[15.1.2]
\label{IV.15.1.2-fr}
Les hypothèses étant celles de (15.1.1), supposons de plus que
$\lambda_z(\mathcal F_y)=1$ (resp. que $\mathcal O_y$ soit régulier et
$\operatorname{long}((\mathcal F_y)_z)=1$). Alors il existe au plus une composante
irréductible $X_i$ de $\operatorname{Supp}(\mathcal F)$ contenant $z$ et telle que
la restriction $X_i\to Y$ de $f$ soit universellement ouverte au point $z$. De
plus, si $z_i$ est le point générique de cette composante, on a
$\lambda_{z_i}(\mathcal F_\eta)=1$ (resp.
$\operatorname{long}((\mathcal F_\eta)_{z_i})=1$).
\end{corollary}

Cela résulte aussitôt de (15.1.1) en notant que l'on a nécessairement, par
définition de $X_i$, $(\mathcal F_\eta)_{z_i}\ne0$ (\textbf{I}, 9.1.13).

\begin{corollary}[15.1.3]
\label{IV.15.1.3-fr}
Soient $Y$ un préschéma irréductible localement noethérien de point générique
$\eta$, $f:X\to Y$ un morphisme localement de type fini, $\mathcal F$ un
$\mathcal O_X$-Module cohérent de support $X$, $x$ un point de $X$, $y=f(x)$. On
suppose que~: $1^\circ$ $x$ n'appartient qu'à une seule composante irréductible
$Z$ de $f^{-1}(y)$~; $2^\circ$
$\mathcal F_y=\mathcal F\otimes_{\mathcal O_Y}\mathbf{k}(y)$ est géométriquement
réduit sur $\mathbf{k}(y)$, autrement dit, si $z$ est le point générique de $Z$,
$\lambda_z(\mathcal F_y)=1$ (resp. $\mathcal O_y$ est régulier et
$\operatorname{long}((\mathcal F_y)_z)=1$). Alors, il existe au plus une
composante irréductible $X'$ de $X$ contenant $x$, telle que
$\dim_x(X'\cap f^{-1}(y))=\dim_x(f^{-1}(y))$ et que la restriction $X'\to Y$ de
$f$ soit universellement ouverte aux points génériques des composantes
irréductibles de $X'\cap f^{-1}(y)$ contenant $x$. De plus, s'il existe une telle
composante $X'$ et si $z'$ est son point générique, on a
$\lambda_{z'}(\mathcal F_\eta)=1$ (resp.
$\operatorname{long}((\mathcal F_\eta)_{z'})=1$).
\end{corollary}

\oldpage[IV-3]{226}
En effet, une composante irréductible $X'$ de $X$ contenant $x$ et telle que
$\dim_x(X'\cap f^{-1}(y))=\dim_x(f^{-1}(y))$ contient nécessairement $Z$,
puisqu'une composante irréductible de $X'\cap f^{-1}(y)$ contenant $x$ doit être
de même dimension que la seule composante irréductible $Z$ de $f^{-1}(y)$
contenant $x$. On peut alors appliquer (15.1.2) à $z$.

\begin{remarks}[15.1.4]
\label{IV.15.1.4-fr}
\begin{enumerate}[label=(\roman*)]
  \item Dans l'énoncé de (15.1.1), on ne peut remplacer «~universellement
  ouvert~» par «~équidimensionnel~», comme le montre l'exemple (14.4.10, (ii)) où
  l'on fait $\mathcal F=\mathcal O_X$~; les fibres de $f$ sont alors des schémas
  artiniens réduits, donc (avec les notations introduites \emph{loc. cit.}) on a
  $\lambda_{x'}(\mathcal F_y)=1$, mais il y a deux composantes irréductibles
  $X_1$, $X_2$ de $X$ contenant $x'$, et le second membre de (15.1.1.1) est donc
  égal à $2$, bien que la restriction de $f$ à chacune des composantes $X_1$,
  $X_2$ soit un morphisme équidimensionnel en tout point. Ce même exemple montre
  aussi que dans (15.1.2) et (15.1.3), on ne peut remplacer l'hypothèse
  «~universellement ouvert~» par «~équidimensionnel~».
  \item Il est vraisemblable que pour la validité de l'inégalité (15.1.1.2), on
  ne peut supprimer l'hypothèse que $\mathcal O_y$ est \emph{régulier}.
\end{enumerate}
\end{remarks}

\subsection{Platitude des morphismes universellement ouverts à fibres géométriquement réduites}
\label{subsection:IV.15.2-fr}

\begin{env}[15.2.1]
\label{IV.15.2.1-fr}
On a vu (2.4.6) qu'un morphisme \emph{plat} et localement de présentation finie
est \emph{universellement ouvert}. On a une réciproque partielle de ce résultat
moyennant des hypothèses supplémentaires~:
\end{env}

\begin{theorem}[15.2.2]
\label{IV.15.2.2-fr}
Soient $Y$ un préschéma localement noethérien, $f:X\to Y$ un morphisme localement
de type fini, $\mathcal F$ un $\mathcal O_X$-Module cohérent, $x$ un point de
$X$, $y=f(x)$. On suppose vérifiées les conditions suivantes~:
\begin{enumerate}[label=(\roman*)]
  \item $\operatorname{Supp}(\mathcal F)=X$.
  \item $f$ est universellement ouvert aux points génériques des composantes
  irréductibles de $f^{-1}(y)$ contenant $x$.
  \item $\mathcal F_y=\mathcal F\otimes_{\mathcal O_Y}\mathbf{k}(y)$ est
  géométriquement réduit au point $x$ (4.6.22).
  \item L'anneau $\mathcal O_y$ est réduit.
\end{enumerate}
Alors $\mathcal F$ est $f$-plat au point $x$.
\end{theorem}

\begin{proof}
Comme $\mathcal O_y$ est réduit et noethérien, nous allons appliquer le critère
valuatif de platitude (11.8.1). Considérons donc un homomorphisme local
$\mathcal O_y\to A'$, où $A'$ est un anneau de valuation discrète~; posons
$Y'=\operatorname{Spec}(A')$, $X'=X\times_Y Y'$, $f'=f_{(Y')}:X'\to Y'$, et
soit $x'$ un point de $X'$ dont les projections dans $X$ et $Y'$ sont
respectivement $x$ et le point fermé $y'$ de $Y'$~; si l'on pose
$\mathcal F'=\mathcal F\otimes_Y Y'$, on a
$\operatorname{Supp}(\mathcal F')=X'$ par (\textbf{I}, 9.1.13)~; tout point
générique $z'$ d'une composante irréductible de $f'^{-1}(y')$ contenant $x'$ a
pour projection dans $f^{-1}(y)$ un point générique $z$ d'une composante
irréductible de $f^{-1}(y)$ contenant $x$, en vertu de (4.2.6)~; donc $f'$ est
universellement ouvert au point $z'$. Enfin, il résulte de (4.7.11) que
$\mathcal F'_{y'}$ est géométriquement réduit au point $x'$. On voit donc qu'on
est ramené à démontrer le

\oldpage[IV-3]{227}
\begin{lemma}[15.2.2.1]
\label{IV.15.2.2.1-fr}
Soient $Y=\operatorname{Spec}(A)$ le spectre d'un anneau de valuation discrète,
$y$ son point fermé, $f:X\to Y$ un morphisme localement de type fini, $x$ un
point de $f^{-1}(y)$, $\mathcal F$ un $\mathcal O_X$-Module cohérent. On suppose
vérifiées les conditions suivantes~:
\begin{enumerate}[label=(\roman*)]
  \item $\operatorname{Supp}(\mathcal F)=X$.
  \item $f$ est ouvert aux points génériques des composantes irréductibles de
  $f^{-1}(y)$ contenant $x$.
  \item $\mathcal F_y=\mathcal F\otimes_{\mathcal O_y}\mathbf{k}(y)$ est réduit
  au point $x$ (3.2.2).
\end{enumerate}
Alors $\mathcal F$ est $f$-plat au point $x$.
\end{lemma}

Désignons en effet par $t$ une uniformisante de $A$, posons $B=\mathcal O_x$ et
$\mathcal F_x=M$~; la condition (i) implique que
$\operatorname{Supp}(M)=\operatorname{Spec}(B)$ (\textbf{I}, 9.1.13). La
condition (ii) implique, en vertu de (14.3.2), que toute composante irréductible
de $X$ contenant $x$ domine $Y$~; en d'autres termes, l'image réciproque dans $A$
de tout idéal premier minimal $\mathfrak p_i$ de $B$ est $0$, ce qui signifie
encore que l'on a $t\notin\mathfrak p_i$ pour tout $i$. Enfin, (iii) signifie que
$M/tM$ est un $(B/tB)$-module \emph{réduit} (3.2.2). Il s'agit de montrer que
sous ces conditions $M$ est un $A$-module \emph{sans torsion}
($\mathbf{0}_{\mathrm I}$, 6.3.4), et pour cela il suffit évidemment de montrer
que $t$ est un élément $M$-\emph{régulier}~; mais cela résulte de (3.4.6.1).
\end{proof}

\begin{corollary}[15.2.3]
\label{IV.15.2.3-fr}
Soient $Y$ un préschéma localement noethérien, $f:X\to Y$ un morphisme localement
de type fini, $x$ un point de $X$, $y=f(x)$. On suppose vérifiées les conditions
suivantes~:
\begin{enumerate}[label=(\roman*)]
  \item $f$ est universellement ouvert aux points génériques des composantes
  irréductibles de $f^{-1}(y)$ contenant $x$.
  \item $f^{-1}(y)$ est géométriquement réduit (sur $\mathbf{k}(y)$) au point
  $x$ (4.6.9).
  \item L'anneau $\mathcal O_y$ est réduit.
\end{enumerate}
Sous ces conditions $f$ est plat au point $x$.
\end{corollary}

Il suffit d'appliquer (15.2.2) à $\mathcal F=\mathcal O_X$ (4.6.22).

\begin{remarks}[15.2.4]
\label{IV.15.2.4-fr}
\begin{enumerate}[label=(\roman*)]
  \item Les trois premières des conditions de (15.2.2) ne changent pas lorsqu'on
  remplace $f$ par $f_{\mathrm{red}}$~; mais si $\mathfrak N$ est le nilradical
  d'un anneau local noethérien $A$, $A/\mathfrak N$ est plat sur $A/\mathfrak N$
  mais non sur $A$ lui-même lorsque $\mathfrak N\ne0$ (Bourbaki,
  \emph{Alg. comm.}, chap.~II, \S~3, n\textsuperscript{o}~2, cor.~2 de la
  prop.~5). On voit donc qu'on ne peut dans (15.2.2) supprimer la condition (iv).

  \item Il résulte de (2.4.6) que si la conclusion de (15.2.2) est vraie, ainsi
  que l'hypothèse (i), $f$ est universellement ouvert dans un voisinage de $x$,
  et en particulier aux points génériques des composantes irréductibles de
  $f^{-1}(y)$ contenant $x$. Même lorsque (i), (iii) et (iv) sont vérifiées, on
  ne peut remplacer (ii) par l'hypothèse plus faible que $f$ est universellement
  ouvert \emph{au point $x$ seulement}. C'est ce que montre l'exemple
  (14.1.3, (i)), où $f$ est évidemment universellement ouvert en tout point de
  $X_2$, donc en particulier au point $x$, intersection de $X_1$ et $X_2$~; les
  conditions (ii) et (iii) de (15.2.3) sont aussi vérifiées par cet exemple.
  Toutefois, $\mathcal O_x$ n'est pas un $\mathcal O_y$-module sans torsion,
  $\mathcal O_y$ s'identifiant à un sous-anneau de $\mathcal O_x$ qui contient
  des diviseurs de $0$ dans $\mathcal O_x$~; donc $f$ n'est pas plat au point $x$.

  \item Dans (15.2.3), la conclusion n'est plus nécessairement valable lorsqu'on
  remplace l'hypothèse (ii) par l'hypothèse plus faible que $f^{-1}(y)$,
  considéré comme $\mathbf{k}(y)$-préschéma, n'a pas de cycles premiers associés
  immergés. Un exemple est fourni en prenant pour $Y$ une courbe ayant un point
  de rebroussement $y$, par exemple
  $Y=\operatorname{Spec}(K[S,T]/(S^3-T^2))$, $K$ étant un corps algébriquement
  clos, et pour $X$ sa normalisée (\textbf{II}, 6.3.8). Si $s$ et $t$ sont les
  classes de $S$ et $T$ dans $A=K[S,T]/(S^3-T^2)$, on vérifie aussitôt que
  $X=\operatorname{Spec}(B)$, où $B=K[s,t,u]$, $u$ étant l'élément $t/s$ du corps
  des fractions de $A$, et comme $s=u^2$, $t=u^3$, on a aussi $B=K[u]$,
  isomorphe à l'anneau de polynômes à une indéterminée sur $K$. Le seul point $x$
  de $X$ au-dessus du point $y$ correspondant à l'idéal maximal $(s)+(t)$
  correspond à l'idéal maximal $(u)$~; mais comme la classe $\bar u$ de $u$ dans
  $\mathcal O_x/\mathfrak m_y\mathcal O_x$ est telle que $\bar u^2=0$,
  $f^{-1}(y)$ n'est pas un $\mathbf{k}(y)$-préschéma réduit. Ici $f$ est un
  morphisme fini, surjectif et radiciel, donc un homéomorphisme universel (2.4.5),
  mais n'est pas plat au point $x$, car si $\mathcal O_x$ était un
  $\mathcal O_y$-module plat, il serait un $\mathcal O_y$-module
  \oldpage[IV-3]{228}
  libre de rang $1$ engendré par l'élément $1$ de $\mathcal O_x$ (Bourbaki,
  \emph{Alg. comm.}, chap.~II, \S~3, n\textsuperscript{o}~2, prop.~5), ce qui est
  absurde.

  \item Montrons enfin que dans (15.2.2) ou (15.2.3), on ne peut remplacer
  «~géométriquement réduit~» par «~réduit~». Nous allons en effet définir deux
  anneaux $A$, $A'$ ayant les propriétés suivantes~:
  \begin{enumerate}[label=\arabic*)]
    \item $A$ est un anneau local noethérien d'idéal maximal $\mathfrak m$,
    intègre, complet, de dimension $1$ et géométriquement unibranche.
    \item $A'$ est la clôture intégrale de $A$, une $A$-algèbre finie dont
    l'idéal maximal est $\mathfrak mA'$, et le corps résiduel $k'$ une extension
    radicielle finie non triviale du corps résiduel $k=A/\mathfrak m$ de $A$.
  \end{enumerate}
  On peut alors prendre $Y=\operatorname{Spec}(A)$,
  $X=\operatorname{Spec}(A')$, $\mathcal F=\mathcal O_X$, $y$ et $x$ étant les
  points fermés de $Y$ et $X$ respectivement~; les hypothèses (i) et (iv) de
  (15.2.2) sont trivialement vérifiées~; comme $A$ est de dimension $1$,
  $f:X\to Y$ est radiciel, fini et surjectif, donc un homéomorphisme universel
  (2.4.5), ce qui prouve l'hypothèse (i) de (15.2.3). Enfin, il est clair que
  $f^{-1}(y)$ est réduit. Cependant $f$ n'est pas plat au point $x$, car si $A'$
  était un $A$-module plat, ce serait un $A$-module libre de rang $1$ (puisqu'il
  a même corps des fractions que $A$) engendré par l'élément $1$ de $A$
  (Bourbaki, \emph{Alg. comm.}, chap.~II, \S~3, n\textsuperscript{o}~2,
  prop.~5), ce qui est absurde.

  Pour construire les anneaux $A$ et $A'$, partons d'un corps imparfait $k$ de
  caractéristique $p>0$~; soient $K_0=k((T))$ le corps des séries formelles sur
  $k$, $A_0$ l'anneau de valuation pour $K_0$ correspondant à la valuation
  discrète égale à l'\emph{ordre} des séries formelles~; le corps résiduel de
  $A_0$ est donc $k$. Soit $\alpha$ un élément de $k$ qui n'est pas puissance
  $p$-ième, et prenons $k'=k(\alpha^{1/p})$~; $K=k'((T))$ est alors une extension
  radicielle finie de $K_0$, et $A'=A_0\otimes_k k'$ l'anneau de valuation
  discrète pour $K$ correspondant à l'ordre des séries formelles sur $K$. Si
  $\mathfrak m'$ est l'idéal maximal de $A'$, on répondra aux conditions 1) et 2)
  en prenant $A=A_0+\mathfrak m'$~: en effet, comme $A_0$ est noethérien et $A'$
  un $A_0$-module de type fini, $A$ est une $A_0$-algèbre finie, donc un anneau
  noethérien~; comme $A'$ est fini sur $A$, $\mathfrak m'$ est le seul idéal
  maximal de $A$, qui est donc un anneau local complet, évidemment de dimension
  $1$ (étant fini sur $A_0$) et géométriquement unibranche puisque $A'$ est
  intégralement clos et a même corps de fractions $K$ que $A$.

  \item La démonstration de (15.2.2) montre toutefois que l'on peut affaiblir la
  condition (iii) en y remplaçant «~géométriquement réduit~» par «~réduit~»
  lorsqu'on peut appliquer le critère valuatif de platitude (11.8.1) en utilisant
  seulement des anneaux de valuation discrète $A'$ dont le corps résiduel $k'$
  est séparable sur $\mathbf{k}(y)$~; en effet, si $x'$ et $z'$ ont les mêmes
  significations que dans (15.2.2), les multiplicités radicielles
  $\mu_i(\mathbf{k}(z')\mid k')$ et
  $\mu_i(\mathbf{k}(z)\mid\mathbf{k}(y))$ sont les mêmes, en vertu de (4.7.3),
  donc il résulte de (4.7.9) que $\operatorname{long}((\mathcal F_y)_z)$ et
  $\operatorname{long}((\mathcal F'_{y'})_{z'})$ sont égales, et l'on est encore
  ramené au cas où $Y$ est le spectre d'un anneau de valuation discrète et $y$
  son point fermé. Par exemple, on pourra obtenir ce résultat plus fort lorsque
  $\mathcal O_y$ est unibranche et dominé par un anneau de valuation discrète
  dont le corps résiduel est extension séparable de $\mathbf{k}(y)$, en vertu de
  (11.6.4). C'est le cas lorsque $\mathcal O_y$ est un anneau régulier, d'après
  (15.1.1.6). Mais, comme nous l'a signalé Hironaka, il y a des anneaux locaux
  intégralement clos noethériens, de dimension $2$ (provenant de schémas
  algébriques sur des corps imparfaits) qui ne satisfont pas à la condition
  précédente.
\end{enumerate}
\end{remarks}

\subsection{Applications : critères de réduction et d'irréductibilité}
\label{subsection:IV.15.3-fr}

\begin{proposition}[15.3.1]
\label{IV.15.3.1-fr}
\oldpage[IV-3]{228}
Soient $Y$ un préschéma localement noethérien, $f:X\to Y$ un morphisme localement
de type fini, $\mathcal F$ un $\mathcal O_X$-Module cohérent, $x$ un point de
$X$, $y=f(x)$. Supposons vérifiées les conditions (i), (ii) et (iii) de
(15.2.2). Alors~:
\begin{enumerate}[label=(\roman*)]
  \item Il existe un voisinage $U$ de $x$ dans $X$ tel que $f|U$ soit
  universellement ouvert et que, pour tout $x'\in U$,
  $\mathcal F_{f(x')}$ soit géométriquement réduit sur
  $\mathbf{k}(f(x'))$ au point $x'$.
  \item Si de plus $Y$ est réduit au point $y$, il existe un voisinage
  $U_1\subset U$ de $x$ tel que $\mathcal F|U_1$ soit $f$-plat et réduit.
\end{enumerate}
\end{proposition}

\begin{proof}
\oldpage[IV-3]{229}
Compte tenu de (\textbf{I}, 5.1.8), (4.2.7) et (4.7.11), les hypothèses ne
changent pas, ni la conclusion (i) de l'énoncé, si l'on remplace $Y$ par
$Y_{\mathrm{red}}$ et $X$ par $X\times_Y Y_{\mathrm{red}}$, et l'on peut donc
supposer $Y$ \emph{réduit}. Il résulte alors de (15.2.2) que $\mathcal F$ est
$f$-plat au point $x$, donc aussi (11.1.1) dans un voisinage de $x$, et
\emph{a fortiori} $f$ est universellement ouvert dans ce voisinage (2.4.6). Le
fait que $\mathcal F_{f(x')}$ soit alors géométriquement réduit au point $x'$
pour tous les points d'un voisinage de $x$ résulte de (12.1.1, (vii)).
L'assertion (ii) résulte donc de ce qui précède et de (3.3.4).
\end{proof}

\begin{proposition}[15.3.2]
\label{IV.15.3.2-fr}
Les notations étant celles de (15.3.1), supposons que les conditions (i), (ii),
(iii) de (15.2.2) soient vérifiées, et de plus que $f^{-1}(y)$ soit
équidimensionnel au point $x$. Alors $f$ est équidimensionnel au point $x$.
\end{proposition}

\begin{proof}
On a, pour tout $y'\in Y$, $\operatorname{Supp}(\mathcal F_{y'})=f^{-1}(y')$.
Par hypothèse $\mathcal F_y$ est réduit (donc vérifie $(S_1)$) et est
équidimensionnel au point $x$. En vertu de (12.1.1, (ii)), on peut donc supposer
que $\mathcal F_{f(x')}$ est équidimensionnel et de dimension constante
$e=\dim_x(f^{-1}(y))$ au point $x'$ pour tous les $x'\in U$, ce qui signifie que
$f^{-1}(f(x'))$ est équidimensionnel et de dimension $e$ au point $x'$. Comme
par ailleurs $\mathcal F|U$ est $f$-plat, il résulte de (2.3.4) et de
(13.3.1, $a$)) que $f$ est équidimensionnel au point $x$.
\end{proof}

\begin{corollary}[15.3.3]
\label{IV.15.3.3-fr}
Les notations étant celles de (15.3.1), supposons vérifiées les conditions
suivantes~:
\begin{enumerate}[label=(\roman*)]
  \item $\operatorname{Supp}(\mathcal F)=X$.
  \item $f$ est universellement ouvert aux points génériques des composantes
  irréductibles de $f^{-1}(y)$ contenant $x$.
  \item $\mathcal F_y$ est géométriquement ponctuellement intègre sur
  $\mathbf{k}(y)$ au point $x$ (4.6.22).
  \item $Y$ est géométriquement unibranche au point $y$.
\end{enumerate}
Alors $x$ n'appartient qu'à une seule composante irréductible de $X$.

Si de plus $Y$ est intègre au point $y$ et $\mathcal F=\mathcal O_X$, alors $X$
est intègre au point $x$ et $f$ est plat au point $x$.
\end{corollary}

\begin{proof}
Notons d'abord que (i) et (iii) entraînent que $x$ n'appartient qu'à une seule
composante irréductible $Z$ de $f^{-1}(y)$. Il résulte donc de (14.4.7) et de
(15.3.2) que pour toute composante irréductible $X_i$ de $X$ contenant $x$
(donc $Z$), la restriction $f_i$ de $f$ à $X_i$ est un morphisme
équidimensionnel au point $x$, et universellement ouvert au point générique de
$Z$. La première assertion de l'énoncé résulte alors de (15.1.3) et la seconde
de (15.3.1).
\end{proof}

\begin{remarks}[15.3.4]
\label{IV.15.3.4-fr}
\begin{enumerate}[label=(\roman*)]
  \item L'exemple (14.4.10, (ii)) montre que, dans (15.3.3), on ne peut
  supprimer l'hypothèse que $Y$ est géométriquement unibranche au point $y$.
  \item La remarque (15.2.4, (vi)) montre que la conclusion de (15.3.3) est
  encore vraie sous les hypothèses suivantes~: $\mathcal O_y$ est régulier,
  $\operatorname{Supp}(\mathcal F)=X$ et $\mathcal F_y$ est intègre au point
  $x$. Nous ignorons si dans cet énoncé, on peut remplacer l'hypothèse que
  $\mathcal O_y$ est régulier par celle qu'il est géométriquement unibranche, ou
  même intègre et intégralement clos.
\end{enumerate}
\end{remarks}

\subsection{Compléments sur les morphismes de Cohen-Macaulay}
\label{subsection:IV.15.4-fr}

\begin{proposition}[15.4.1]
\label{IV.15.4.1-fr}
Soient $Y$ un préschéma localement noethérien, $f:X\to Y$ un morphisme localement
de type fini, $\mathcal F$ un $\mathcal O_X$-Module cohérent tel que
$\operatorname{Supp}(\mathcal F)=X$, $x$ un point de $X$, $y=f(x)$. On pose,
pour tout $z\in X$, $X_{f(z)}=f^{-1}(f(z))$,
$\mathcal F_{f(z)}=\mathcal F\otimes_{\mathcal O_{f(z)}}\mathbf{k}(f(z))$.
Supposons que $\mathcal F$ soit $f$-plat au point $x$ et que $\mathcal F_y$ soit
un $\mathcal O_{X_y}$-Module de Cohen-Macaulay au point $x$. Alors $f$ est
\oldpage[IV-3]{230}
équidimensionnel au point $x$.

En effet ((11.1.1) et (12.1.1, (vi))), il existe un voisinage $U$ de $x$ tel que
$\mathcal F|U$ soit $f$-plat et que pour tout $x'\in U$,
$\mathcal F_{f(x')}$ soit un $\mathcal O_{X_{f(x')}}$-Module de Cohen-Macaulay
au point $x'$~; cette dernière propriété entraîne que $\mathcal F_{f(x')}$ est
équidimensionnel au point $x'$ et que $x'$ n'appartient à aucun cycle premier
immergé associé à $\mathcal F_{f(x')}$ ($\mathbf{0}_{\mathrm{IV}}$, 16.5.4). En
outre, en vertu de (12.1.1, (ii)), on peut supposer que les dimensions des
composantes irréductibles de
$\operatorname{Supp}(\mathcal F_{f(x')})|(U\cap X_{f(x')})$ aient une (même)
valeur indépendante de $x'$. Comme
$\operatorname{Supp}(\mathcal F_{f(x')})=X_{f(x')}$ par hypothèse, il résulte de
(2.3.4) et de (13.3.1, $a$)) que $f$ est équidimensionnel au point $x$.
\end{proposition}

\begin{proposition}[15.4.2]
\label{IV.15.4.2-fr}
Soient $Y$ un préschéma localement noethérien, $f:X\to Y$ un morphisme localement
de type fini, $\mathcal F$ un $\mathcal O_X$-Module cohérent tel que
$\operatorname{Supp}(\mathcal F)=X$, $x$ un point de $X$, $y=f(x)$. On suppose
que $\mathcal O_y$ est régulier et que $\mathcal F$ est un $\mathcal O_X$-Module
de Cohen-Macaulay au point $x$. Alors (avec les notations de (15.4.1)) les
conditions suivantes sont équivalentes~:
\begin{enumerate}[label=\emph{\alph*})]
  \item $\mathcal F$ est $f$-plat au point $x$ et $\mathcal F_y$ est un
  $\mathcal O_{X_y}$-Module de Cohen-Macaulay au point $x$.
  \item $\mathcal F$ est $f$-plat au point $x$.
  \item $f$ est universellement ouvert dans un voisinage de $x$.
  \item $f$ est ouvert aux points génériques des composantes irréductibles de
  $X_y$ contenant $x$.
  \item $\dim(\mathcal F_x)=\dim(\mathcal O_y)+\dim((\mathcal F_y)_x)$.
  \item[\emph{e}$'$)]
  $\dim(\mathcal O_{X,x})=\dim(\mathcal O_y)+
  \dim(\mathcal O_{X,x}\otimes_{\mathcal O_y}\mathbf{k}(y))$.
\end{enumerate}
Comme $\operatorname{Supp}(\mathcal F)=X$, les conditions $e)$ et $e')$ sont
équivalentes par définition (5.1.12). La condition $a)$ implique trivialement
$b)$~; $b)$ entraîne que $\mathcal F$ est $f$-plat aux points d'un voisinage de
$x$ (11.1.1), donc $b)$ entraîne $c)$ (2.4.6)~; $c)$ entraîne trivialement $d)$,
et $d)$ entraîne $e')$ par (14.2.1). Enfin, puisque $\mathcal O_y$ est régulier
et que $\mathcal F_x$ est un $\mathcal O_{X,x}$-module de Cohen-Macaulay, il
résulte de (6.1.4) que $e)$ entraîne $a)$.
\end{proposition}

\begin{proposition}[15.4.3]
\label{IV.15.4.3-fr}
Soient $f:X\to Y$ un morphisme localement de présentation finie, $\mathcal F$ un
$\mathcal O_X$-Module de présentation finie, tel que
$\operatorname{Supp}(\mathcal F)=X$, $x$ un point de $X$, $y=f(x)$. Alors (avec
les notations de (15.4.1)) les conditions suivantes sont équivalentes~:
\begin{enumerate}[label=\emph{\alph*})]
  \item $\mathcal F$ est $f$-plat au point $x$ et $\mathcal F_y$ est un
  $\mathcal O_{X_y}$-Module de Cohen-Macaulay au point $x$.
  \item Il existe un voisinage ouvert $U$ de $x$ dans $X$, et un $Y$-morphisme
  quasi-fini $g:U\to Y'=Y[T_1,\ldots,T_n]$ ($T_i$ indéterminées), tels que
  $\mathcal F|U$ soit $g$-plat au point $x$.
\end{enumerate}
\end{proposition}

\begin{proof}
Montrons d'abord que $b)$ entraîne $a)$. Posons $h:Y'\to Y$~; le
$\mathbf{k}(y)$-préschéma $Y_y'=h^{-1}(y)$ est régulier
($\mathbf{0}_{\mathrm{IV}}$, 17.3.7)~; soient $A$ l'anneau local de ce préschéma
au point $y'=g(x)$, $k$ son corps résiduel, et soit $B$ l'anneau local de $X_y$
au point $x$~; posons d'autre part $(\mathcal F_y)_x=M$, qui est un $B$-module
de type fini. L'hypothèse que le morphisme $g$ est quasi-fini entraîne qu'il en
est de même de $g_y:X_y\to Y_y'$ (\textbf{II}, 6.2.4), donc $M\otimes_A k$ est
un $(B\otimes_A k)$-module de longueur finie (\textbf{II}, 6.2.2)~; comme par
hypothèse $M$ est un $A$-module plat et $A$ un anneau de Cohen-Macaulay, $M$ est
un $B$-module de Cohen-Macaulay (6.3.3). Le fait que $\mathcal F$ est $f$-plat
au point $x$ résulte de ce qu'il est $g$-plat et de ce que le morphisme $h$ est
plat (2.1.6).

Montrons maintenant que $a)$ entraîne $b)$. La question étant locale sur $X$ et
$Y$, on peut supposer $X$ et $Y$ affines~; utilisant (8.9.1) et (11.2.7), on se
ramène au cas où $X$ et $Y$ sont noethériens, compte tenu de (6.7.1).
L'hypothèse entraîne que $f$ est
\oldpage[IV-3]{231}
équidimensionnel au point $x$ (15.4.1), d'où l'existence d'un morphisme
quasi-fini $g:U\to Y'$ tel que toute composante irréductible de $U$ domine $Y'$
(13.3.1, $b$)) et que $g_y:U_y\to Y_y'$ soit fini et surjectif (13.3.1.1).
D'autre part, pour prouver que $\mathcal F|U$ est $g$-plat au point $x$, il
suffit, puisque $h:Y'\to Y$ est plat, de montrer (en vertu du critère de
platitude par fibres (11.3.10)) que $\mathcal F_y$ est $g_y$-plat au point $x$.
Mais par hypothèse $\mathcal F_y$ est un $\mathcal O_{X_y}$-Module de
Cohen-Macaulay au point $x$ et $Y_y'$ un préschéma \emph{régulier}~; d'autre
part, comme $g_y$ est fini et surjectif, il vérifie la condition $e')$ de
(15.4.2) en vertu de (5.6.10)~; il suffit donc pour conclure d'appliquer (15.4.2)
à $g_y$ et à $\mathcal F_y$.
\end{proof}

% IV3_B32_P231_END_BEFORE_SUBSECTION_15.5

\subsection{Rang séparable des fibres d'un morphisme quasi-fini et universellement ouvert. Application aux composantes connexes géométriques des fibres d'un morphisme propre}
\label{subsection:IV.15.5-fr}

\begin{proposition}[15.5.1]
\label{IV.15.5.1-fr}
Soient $Y$ un préschéma localement noethérien, $f:X\to Y$ un morphisme séparé et quasi-fini. Pour tout $z\in Y$, soit $n(z)$ le nombre géométrique des points de $f^{-1}(z)$ (ou rang séparable de $f^{-1}(z)$ sur $\mathbf{k}(z)$, cf. \textbf{I}, 6.4.8). On a les propriétés suivantes~:
\begin{enumerate}[label=\emph{(\roman*)}]
  \item La fonction $z\mapsto n(z)$ est semi-continue inférieurement en tout point $y\in Y$ tel que $f$ soit universellement ouvert aux points de $f^{-1}(y)$.
  \item Soit $y$ un point de $Y$ tel que $f$ soit universellement ouvert aux points de $f^{-1}(y)$~; si la fonction $z\mapsto n(z)$ est constante dans un voisinage de $y$, alors $f$ est propre au point $y$ (15.7.1).
  \item On suppose que $f$ soit universellement ouvert et que tout point de $Y$ soit géométriquement unibranche~; alors, si la fonction $z\mapsto n(z)$ est localement constante dans $Y$, les composantes irréductibles de $X$ sont deux à deux disjointes.
\end{enumerate}
\end{proposition}

\begin{proof}
(i) Notons que puisque $f$ est quasi-fini, le nombre $n(z)$ est le nombre géométrique des composantes connexes de $f^{-1}(z)$~; on sait que la fonction $z\mapsto n(z)$ est localement constructible (9.7.9)~; compte tenu de ($\mathbf{0}_{\mathrm{III}}$, 9.3.4), on est ramené à prouver le

\begin{corollary}[15.5.2]
\label{IV.15.5.2-fr}
Sous les hypothèses générales de (15.5.1), soit $y$ un point de $Y$ tel que $f$ soit universellement ouvert aux points de $f^{-1}(y)$. Alors, si $y'$ est une générisation de $y$, on a
\begin{equation}
\tag{15.5.2.1}
\label{IV.15.5.2.1-fr}
  n(y)\leqslant n(y').
\end{equation}
\end{corollary}

Pour tout changement de base $g:Z\to Y$, $f_{(Z)}:X_{(Z)}\to Z$ est séparé, quasi-fini et universellement ouvert aux points de $X_{(Z)}$ se projetant dans $f^{-1}(y)$~; en outre, si $z,z'$ sont deux points de $Z$ tels que $g(z)=y$, $g(z')=y'$ et que $z'$ soit une générisation de $z$, on a $n(z)=n(y)$ et $n(z')=n(y')$ (I, 6.4.12)~; il suffit donc de démontrer (15.5.2.1) pour un $g$ convenable. On peut évidemment supposer $y\ne y'$. Nous utiliserons le lemme suivant~:

\begin{lemma}[15.5.2.2]
\label{IV.15.5.2.2-fr}
Sous les hypothèses générales de (15.5.1), si $y$ est un point de $Y$, $y'$ une générisation de $y$ distincte de $y$, il existe un spectre d'anneau de valuation discrète $Z$, de point fermé $z$ et de point générique $z'$, et un morphisme $g:Z\to Y$ tels que~:
\begin{enumerate}[label=\emph{\arabic*\textsuperscript{\textdegree}}]
  \item $g(z)=y$, $g(z')=y'$~;
  \item si l'on pose $f'=f_{(Z)}$, $n(y)$ est le nombre de points de $f'^{-1}(z)$ et $n(y')$ est le nombre de points de $f'^{-1}(z')$.
\end{enumerate}
\end{lemma}

En effet, il existe un anneau de valuation discrète $A_1$ et un morphisme $g_1$ de $Z_1=\operatorname{Spec}(A_1)$ dans $Y$ tels que, si $z_1$ et $z'_1$ sont le point fermé et le point générique de $Z_1$,
\oldpage[IV-3]{232}
on ait $g_1(z_1)=y$ et $g_1(z'_1)=y'$ (\textbf{II}, 7.1.9)~; on peut donc déjà supposer que $Y$ est le spectre d'un anneau de valuation discrète $A$, $y$ son point fermé et $y'$ son point générique. Pour tout $x\in f^{-1}(y)$, $\mathbf{k}(x)$ est une extension finie de $\mathbf{k}(y)$ par hypothèse (\textbf{II}, 6.2.2)~; on déduit de (4.5.11) qu'il existe une extension algébrique finie $k$ de $\mathbf{k}(y)$ telle que le nombre de points de $f^{-1}(y)\otimes_{\mathbf{k}(y)}k$ soit égal à $n(y)$. Comme il y a un anneau de valuation discrète $A_2$ dominant $A$ et dont le corps résiduel est fini sur $\mathbf{k}(y)$ et contient $k$ (\textbf{II}, 7.1.2), on peut en second lieu supposer que $n(y)$ est le nombre de points de $f^{-1}(y)$. Enfin, le même raisonnement montre qu'il y a une extension finie $K$ de $\mathbf{k}(y')$ telle que le nombre de points de $f^{-1}(y')\otimes_{\mathbf{k}(y')}K$ soit égal à $n(y')$~; si $w$ est une valuation de $\mathbf{k}(y')$ associée à $A$, il y a une valuation discrète de $K$ prolongeant $w$, et l'anneau $A_3$ de cette valuation domine $A$, a $K$ pour corps des fractions, et répond aux conditions du lemme.

On peut donc désormais supposer que $Y$ est le spectre d'un anneau de valuation discrète, $y$ son point fermé, $y'$ son point générique, et que $n(y)$ et $n(y')$ sont les nombres de points de $f^{-1}(y)$ et $f^{-1}(y')$, de sorte que pour tout $x\in f^{-1}(y)$ (resp. tout $x'\in f^{-1}(y')$) on a $\mathbf{k}(x)=\mathbf{k}(y)$ (resp. $\mathbf{k}(x')=\mathbf{k}(y')$).

Désignons alors par $X_i$ les sous-préschémas réduits de $X$ ayant pour espaces sous-jacents les composantes irréductibles de $X$ ($1\leqslant i\leqslant m$). En vertu de (1.10.4), la restriction $X_i\to Y$ de $f$ à tout $X_i$ est un morphisme dominant, et comme $f^{-1}(y')$ est discret, chacune des intersections $X_i\cap f^{-1}(y')$ est réduite au point générique de $X_i$~; d'où (en désignant par $n_i(z)$ le nombre géométrique de points de $X_i\cap f^{-1}(z)$ pour tout $z\in Y$), $n_i(y')=1$ pour tout $i$, et $n(y')=\sum_{i=1}^{m}n_i(y')=m$. Par ailleurs, $f^{-1}(y)$ est réunion des $X_i\cap f^{-1}(y)$, d'où
\begin{equation}
\tag{15.5.2.3}
\label{IV.15.5.2.3-fr}
  n(y)\leqslant\sum_{i=1}^{m}n_i(y),
\end{equation}
et, pour prouver (15.5.2.1), il suffit d'établir que $n_i(y)\leqslant 1$ pour tout $i$. Autrement dit, on peut supposer $X$ intègre, et le morphisme $f$ birationnel. Mais alors, pour tout $x\in f^{-1}(y)$, on a $\mathcal O_y\subset\mathcal O_x\subset K$, où $K=\mathbf{k}(y')$ est le corps des fractions de $\mathcal O_y$. Comme $\mathcal O_y$ est un anneau de valuation et que $\mathcal O_x$ domine $\mathcal O_y$, on a nécessairement $\mathcal O_y=\mathcal O_x$ (\textbf{II}, 7.1.1). Il existe alors une $Y$-section $h:Y\to X$ telle que $h(y)=x$ (\textbf{I}, 2.4.4)~; comme $Y$ est réduit et que $X$ est séparé sur $Y$, une telle section est unique (\textbf{I}, 7.2.2), donc l'ensemble $f^{-1}(y)$ ne contient qu'un seul point, ce qui achève de prouver (i).

(ii) Comme au début de (i), on est ramené à montrer que si les deux membres de (15.5.2.1) sont égaux pour toute générisation $y'$ de $y$, $f$ est propre au point $y$. Grâce au critère de propreté locale (15.7.5) et aux remarques du début on peut encore supposer que $Y=\operatorname{Spec}(A)$ est le spectre d'un anneau de valuation discrète $A$, $y$ son point fermé et $y'$ son point générique, et il s'agit alors de montrer que $f$ est propre. Utilisons (15.5.2.2) et notons que si $A'$ est un anneau de valuation discrète dominant $A$, $A'$ est un $A$-module sans torsion, donc le morphisme $\operatorname{Spec}(A')\to\operatorname{Spec}(A)$ est fidèlement plat et quasi-compact, et il revient donc au même de dire que $f$ est propre ou que le morphisme déduit de $f$ par le changement de base $\operatorname{Spec}(A')\to\operatorname{Spec}(A)$ est
\oldpage[IV-3]{233}
propre (2.7.1, (vii))~; on peut donc supposer que $n(y)$ et $n(y')$ sont les nombres de points des fibres $f^{-1}(y)$ et $f'^{-1}(y')$ respectivement. Avec les notations de (i), on a alors par (15.5.2.3) et (15.5.2.1)
\begin{equation}
\tag{15.5.2.4}
\label{IV.15.5.2.4-fr}
  n(y)\leqslant\sum_i n_i(y)\leqslant\sum_i n_i(y')=n(y').
\end{equation}
et, pour que les membres extrêmes soient égaux, il faut que $n_i(y)=n_i(y')=1$ pour tout $i$. Comme on l'a vu dans (i), si $n_i(y)=1$ pour tout $i$, $X_i\cap f^{-1}(y)$ n'est pas vide et la restriction $X_i\to Y$ de $f$ est un isomorphisme~; comme les $X_i$ sont des sous-préschémas fermés de $X$, cela entraîne que $f$ est propre (\textbf{II}, 5.4.5).

(iii) On peut se borner au cas où $Y$ est affine et réduit, et comme $Y$ est par hypothèse localement intègre (\textbf{I}, 5.1.4), on peut supposer $Y$ intègre, et la fonction $y\mapsto n(y)$ constante dans $Y$. Soient encore $X_i$ ($1\leqslant i\leqslant m$) les composantes irréductibles de $X$. Il résulte de (1.10.4) que si $y'$ est le point générique de $Y$, chacun des $X_i\cap f^{-1}(y')$ est réduit à un seul point~; la restriction $f_i:X_i\to Y$ de $f$ étant un morphisme quasi-fini et dominant et tout point de $Y$ étant géométriquement unibranche, il résulte du critère de Chevalley (14.4.4) que $f_i$ est universellement ouvert. Si alors $y$ est un point quelconque de $Y$, on a $n_i(y)\leqslant n_i(y')$~; d'autre part, comme les $X_i\cap f^{-1}(y)$ sont deux à deux disjoints, on a $\sum_i n_i(y')=n(y')$~; on a donc encore les relations (15.5.2.4). Mais par hypothèse les membres extrêmes sont égaux, donc $n(y)=\sum_i n_i(y)$, ce qui implique que les $X_i\cap f^{-1}(y)$ sont deux à deux disjoints, et achève la démonstration.
\end{proof}

Combinant cette proposition avec le théorème de connexion de Zariski et ses conséquences (\textbf{III}, 4.3), on obtient les résultats suivants qui complètent (\textbf{III}, 4.3.7 et 4.3.10)~:

\begin{proposition}[15.5.3]
\label{IV.15.5.3-fr}
Soient $Y$ un préschéma irréductible localement noethérien, $f:X\to Y$ un morphisme propre~; pour tout $y\in Y$, soit $n(y)$ le nombre géométrique de composantes connexes (4.5.2) de $f^{-1}(y)$. On suppose que la restriction de $f$ à toute composante irréductible de $X$ soit un morphisme dominant dans $Y$. Alors, si $y_0$ est un point géométriquement unibranche de $Y$, la fonction $y\mapsto n(y)$ est semi-continue inférieurement au point $y_0$.
\end{proposition}

\begin{proof}
Considérons la factorisation de Stein $f=g\circ f'$ de $f$ (\textbf{III}, 4.3.3), où $g:Y'\to Y$ est un morphisme fini et $f':X\to Y$ est un morphisme surjectif dont les fibres sont géométriquement connexes (\textbf{III}, 4.3.4). Puisque $f'$ est surjectif, l'image réciproque par $f'$ de toute composante irréductible de $Y'$ contient une composante irréductible de $X$ au moins~; donc chacune des composantes irréductibles de $Y'$ domine $Y$. On peut donc appliquer le critère de Chevalley (14.4.4) aux restrictions de $g$ à chacune de ces composantes irréductibles, et on en conclut que $g$ est universellement ouvert en tout point de $g^{-1}(y_0)$. La conclusion résulte donc de (15.5.2) et du fait que le nombre géométrique des composantes connexes de $f^{-1}(y)$ est égal au nombre géométrique des points de $g^{-1}(y)$ (\textbf{III}, 4.3.4).
\end{proof}

\begin{corollary}[15.5.4]
\label{IV.15.5.4-fr}
Soient $Y$ un préschéma localement noethérien, $f:X\to Y$ un morphisme propre~; soit $y$ un point de $Y$ tel que $f$ soit universellement ouvert aux points de $f^{-1}(y)$~;
\oldpage[IV-3]{234}
alors, avec les notations de (15.5.3), la fonction $z\mapsto n(z)$ est semi-continue inférieurement au point $y$.
\end{corollary}

\begin{proof}
Avec les notations de (15.5.3), notons que dans la factorisation de Stein $f=g\circ f'$, $f'$ est surjectif~; comme $f$ est universellement ouvert aux points de $f^{-1}(y)$, $g$ est universellement ouvert aux points de $g^{-1}(y)$ (14.3.4, (i))~; il suffit alors d'appliquer à $g$ la proposition (15.5.1), (i) en tenant compte de (\textbf{III}, 4.3.4).
\end{proof}

\begin{remarks}[15.5.5]
\label{IV.15.5.5-fr}
\begin{enumerate}[label=\emph{(\roman*)}]
  \item Même si $Y$ est intègre et normal, $X$ intègre, $f$ fini et surjectif (donc universellement ouvert en vertu du critère de Chevalley (14.4.4)), la fonction $y\mapsto n(y)$ n'est pas nécessairement localement constante. On en a un exemple en prenant $Y=\operatorname{Spec}(k[T])$ où $k$ est algébriquement clos, et $X=\operatorname{Spec}(A)$, où $A=k[S,T]/(S^2+T^2-1)$~; aux points $y',y''$ de $Y$ correspondant aux idéaux maximaux $(T-1)$ et $(T+1)$, on a $n(y')=n(y'')=1$, mais $n(y)=2$ en tous les autres points de $Y$. On va donner ci-dessous une condition supplémentaire assurant que $y\mapsto n(y)$ est localement constante (15.5.7).
  \item L'exemple (14.4.10, (ii)) montre que dans (15.5.1, (iii)), on ne peut supprimer l'hypothèse que les points de $Y$ sont géométriquement unibranches.
  \item L'exemple (11.7.5) montre que la conclusion de (15.5.1, (i)) n'est plus valable si on suppose seulement que le morphisme $f$ est ouvert (même si, comme c'est le cas dans l'exemple cité, $f$ est fini et surjectif).
  \item Enfin, dans (15.5.1), on ne peut se dispenser de l'hypothèse que le morphisme $f$ est séparé, comme le montre l'exemple où $X$ est la droite affine dont on a «~dédoublé un point~» (\textbf{I}, 5.5.11), et $Y$ la droite affine~: ici $f$ est un isomorphisme local (donc est plat) et est quasi-fini, mais $z\mapsto n(z)$ n'est pas semi-continue inférieurement.
\end{enumerate}
\end{remarks}

\begin{lemma}[15.5.6]
\label{IV.15.5.6-fr}
Soit $Y$ le spectre d'un anneau de valuation discrète, $a$ son point fermé, $b$ son point générique. Soit $f:X\to Y$ un morphisme localement de type fini et ouvert. Supposons que $X$ soit connexe et que la fibre $f^{-1}(a)$ soit un préschéma réduit. Alors $f^{-1}(b)$ est connexe.
\end{lemma}

\begin{proof}
Nous utiliserons le lemme purement topologique suivant (cas particulier d'un résultat plus général du chap. III, 3\textsuperscript{e} Partie)~:
\begin{lemma}[15.5.6.1]
\label{IV.15.5.6.1-fr}
Soit $X$ un espace localement noethérien, dont toute partie fermée irréductible admet un point générique. Soit $Z$ une partie fermée rare de $X$ ayant la propriété suivante~: pour tout $x\in Z$, si l'on désigne par $T_x$ l'ensemble des générisations de $x$ dans $X$, alors $T_x-\{x\}$ est connexe. Dans ces conditions, si $X$ est connexe, $X-Z$ est connexe.
\end{lemma}

\begin{proof}
Donnons-en une démonstration indépendante. Raisonnant par l'absurde, supposons que l'ouvert $U=X-Z$, partout dense dans $X$, soit réunion de deux ouverts non vides sans point commun $U'$ et $U''$, et désignons par $X'$ et $X''$ les adhérences dans $X$ de $U'$ et $U''$. Comme $X=\overline{U}=X'\cup X''$ et que $X$ est connexe, on a $X'\cap X''\ne\varnothing$, et il est clair que $X'\cap X''\subset Z$. Soit $x$ un point générique d'une composante irréductible de $X'\cap X''$. Comme $X$ est localement noethérien, il y a un voisinage ouvert $V$ de $x$ dans $X$ tel que $V\cap U'$ et $V\cap U''$ n'aient qu'un nombre fini de composantes irréductibles~; comme $x$ est adhérent à $U'$ et à $U''$, il est nécessairement dans l'adhérence d'une composante irréductible $Z'$ de $V\cap U'$ et dans l'adhérence d'une composante irréductible $Z''$ de $V\cap U''$. Mais $Z'$ (resp. $Z''$) est fermé et irréductible dans $X$, donc admet un point générique $z'$ (resp. $z''$), et on a nécessairement $z'\in U'$ et $z''\in U''$. Ceci prouve que les intersections de $T_x-\{x\}$ avec $U'$ et $U''$ sont non vides. Mais par définition $(X'\cap X'')\cap(T_x-\{x\})=\varnothing$~; donc les intersections de $X'$ et de $X''$ avec $T_x-\{x\}$ sont deux fermés non vides disjoints dans $T_x-\{x\}$, dont la réunion est $T_x-\{x\}$~; or cela contredit l'hypothèse que $T_x-\{x\}$ est connexe. C.Q.F.D.
\end{proof}
\oldpage[IV-3]{235}
Pour prouver (15.5.6), nous appliquerons le lemme (15.5.6.1) à $X$ et à sa partie fermée $Z=f^{-1}(a)$ qui est rare puisque $f$ est ouvert. Notons que l'hypothèse, compte tenu de (15.2.2.1), implique que $f$ est plat aux points de $f^{-1}(a)$. Pour un tel point $x$, $\mathcal O_x$ est donc un $\mathcal O_a$-module sans torsion ($\mathbf{0}_{\mathrm I}$, 6.3.4), et en particulier, si $t$ est une uniformisante de $\mathcal O_a$, $t$ est $\mathcal O_x$-régulier. En outre $\mathcal O_x/t\mathcal O_x$ est réduit par hypothèse~; s'il est de profondeur $0$, c'est donc un corps ($\mathbf{0}$, 16.4.7), et comme $t\mathcal O_x$ est alors l'idéal maximal de $\mathcal O_x$, $\mathcal O_x$ est un anneau de valuation discrète ($\mathbf{0}$, 17.1.4), donc $\operatorname{Spec}(\mathcal O_x)-\{x\}$ est réduit à un point, et a fortiori est connexe. Si au contraire $\operatorname{prof}(\mathcal O_x/t\mathcal O_x)\geqslant 1$, on a $\operatorname{prof}(\mathcal O_x)\geqslant 2$ puisque $t$ est $\mathcal O_x$-régulier ($\mathbf{0}$, 16.4.6)~; il résulte alors du théorème de Hartshorne (5.10.7) que $\operatorname{Spec}(\mathcal O_x)-\{x\}$ est connexe, ce qui termine la démonstration.
\end{proof}

\begin{proposition}[15.5.7]
\label{IV.15.5.7-fr}
Soient $Y$ un préschéma localement noethérien, $f:X\to Y$ un morphisme propre. Pour tout $z\in Y$, soit $n(z)$ le nombre géométrique de composantes connexes de $f^{-1}(z)$. Soit $y$ un point de $Y$ tel que $f$ soit universellement ouvert aux points de $f^{-1}(y)$ et que $f^{-1}(y)$ soit géométriquement réduit sur $\mathbf{k}(y)$ (4.6.2). Alors la fonction $z\mapsto n(z)$ est constante dans un voisinage de $y$.
\end{proposition}

\begin{proof}
On sait déjà (15.5.4) que sous les conditions de l'énoncé, $z\mapsto n(z)$ est semi-continue inférieurement au point $y$~; tout revient à voir qu'elle est aussi semi-continue supérieurement, et par le même raisonnement qu'au début de la démonstration de (15.5.1, (i)), il suffit de montrer que si $y'$ est une générisation de $y$, on a
\begin{equation}
\tag{15.5.7.1}
\label{IV.15.5.7.1-fr}
  n(y')\leqslant n(y).
\end{equation}

Utilisant le raisonnement du début de la démonstration de (15.5.2) et le lemme (15.5.2.2) appliqué au morphisme fini $g$ de la factorisation de Stein $f=g\circ f'$ de $f$ (en se souvenant que le nombre géométrique de points de $g^{-1}(y)$ est égal à celui des composantes connexes de $f^{-1}(y)$ ($\mathbf{III}$, 4.3.4)), on se ramène au cas où $n(y)$ et $n(y')$ sont respectivement le nombre de composantes connexes de $f^{-1}(y)$ et $f^{-1}(y')$, et où $Y$ est le spectre d'un anneau de valuation discrète, $y$ le point fermé et $y'$ le point générique de $Y$. On peut en outre remplacer $X$ par une quelconque de ses composantes connexes, ces dernières étant ouvertes et fermées dans $X$, et propres sur $Y$. Supposons donc $X$ connexe~; l'hypothèse que $f$ est ouvert aux points de $f^{-1}(y)$ entraîne que si $f^{-1}(y)\ne\varnothing$, on a aussi $f^{-1}(y')\ne\varnothing$ ($\mathbf{0}_{\mathrm I}$, 1.10.3)~; l'hypothèse que $f$ est fermé entraîne que si $f^{-1}(y')\ne\varnothing$, on a aussi $f^{-1}(y)\ne\varnothing$ (\textbf{II}, 7.2.1). Donc, si $X\ne\varnothing$ (ce qu'on peut évidemment supposer, la proposition étant triviale dans le cas contraire), $f^{-1}(y)$ et $f^{-1}(y')$ sont tous deux non vides. En outre, l'hypothèse entraîne que le préschéma $f^{-1}(y)$ est réduit (4.6.1)~; comme $f$ est ouvert aux points de $f^{-1}(y)$, le lemme (15.5.6) implique que $f^{-1}(y')$ est connexe, autrement dit $n(y')=1$. Comme $f^{-1}(y)$ n'est pas vide, cela prouve (15.5.7.1).
\end{proof}

\begin{remark}[15.5.8]
\label{IV.15.5.8-fr}
La relation (15.5.7.1) n'est plus nécessairement valable si $f$ n'est pas supposé propre, même s'il vérifie les autres hypothèses de (15.5.7). Avec les notations de (15.5.5, (i)), il suffit pour le voir de considérer la restriction de $f$ à $X-\{x'\}$, où $x'$ est un des deux points de $X$ au-dessus d'un point $y$ distinct de $y'$ et $y''$~; on a alors $n(y)=1$, tandis que si $\eta$ est le point générique de $Y$, $n(\eta)=2$.
\end{remark}

\oldpage[IV-3]{236}
\begin{proposition}[15.5.9]
\label{IV.15.5.9-fr}
Soit $f:X\to Y$ un morphisme plat et localement de présentation finie~; pour tout $y\in Y$, soit $n(y)$ le nombre géométrique de composantes connexes de $f^{-1}(y)$.
\begin{enumerate}[label=\emph{(\roman*)}]
  \item Si $f$ est séparé et quasi-fini, la fonction $y\mapsto n(y)$ (égale alors au nombre géométrique des points de $f^{-1}(y)$) est semi-continue inférieurement dans $Y$~; si elle est constante dans un voisinage de $y_0$, $f$ est propre au point $y_0$.
  \item Si $f$ est propre, la fonction $y\mapsto n(y)$ est semi-continue inférieurement dans $Y$~; si de plus $f^{-1}(y_0)$ est géométriquement réduit sur $\mathbf{k}(y_0)$, $n(y)$ est constante dans un voisinage de $y_0$.
\end{enumerate}
\end{proposition}

\begin{proof}
Dans tous les cas, les questions sont locales sur $Y$, donc on peut supposer $Y$ affine et $f$ de présentation finie~; utilisant (8.9.1), (8.10.5) et (11.2.7), on se ramène au cas où $Y$ est noethérien (en utilisant en outre (8.2.11) pour la seconde assertion de (i)). Comme de plus $f$ est alors universellement ouvert (2.4.6), il suffit d'appliquer (15.5.1), (15.5.4) et (15.5.7) pour conclure.
\end{proof}

\begin{remark}[15.5.10]
\label{IV.15.5.10-fr}
On notera qu'on a ainsi obtenu une autre démonstration de (12.2.4, (vi)). Inversement, on peut déduire (15.5.7) de (12.2.4, (vi))~: en effet, on peut se borner au cas où $Y$ est réduit (en remplaçant $f$ par $f_{\mathrm{red}}$), et il résulte alors des hypothèses de (15.5.7) et de (15.2.3) que $f$ est plat dans un voisinage ouvert de $f^{-1}(y)$~; comme de plus $f$ est propre, on peut supposer que ce voisinage est de la forme $f^{-1}(U)$, où $U$ est un voisinage ouvert de $y$ dans $Y$. On conclut alors par (12.2.4, (vi)). La démonstration de (15.5.7) donnée plus haut a l'avantage de mettre en évidence le résultat (15.5.6), qui a un intérêt indépendant.
\end{remark}

% IV3_B33_P236_END_BEFORE_SUBSECTION_15.6

\subsection{Composantes connexes des fibres le long d'une section}
\label{subsection:IV.15.6-fr}

\begin{proposition}[15.6.1]
\label{IV.15.6.1-fr}
Soient $Y$ un préschéma localement noethérien, $f:X\to Y$ un morphisme localement de type fini, $g:Y\to X$ une $Y$-section de $X$ (\textbf{I}, 2.5.5). Pour tout $y\in Y$, désignons par $X_y^0$ la composante connexe de $f^{-1}(y)$ contenant le point $g(y)$. Soient $y$ un point de $Y$, $y'$ une générisation de $y$ dans $Y$. Alors :
\begin{enumerate}[label=\emph{(\roman*)}]
  \item S'il existe une composante irréductible $Z$ de $X_{y'}^0$, de dimension égale à $\dim(X_{y'}^0)$ et contenant $g(y')$ (ce qui sera le cas si $X_{y'}^0$ est irréductible), on a
\begin{equation}
\tag{15.6.1.1}
\label{IV.15.6.1.1-fr}
  \dim(X_y^0)\geqslant\dim(X_{y'}^0).
\end{equation}
  \item Si de plus $X_y^0$ est irréductible et si les deux membres de (15.6.1.1) sont égaux, alors $X_y^0\subset\overline{X_{y'}^0}$ (adhérence dans $X$).
\end{enumerate}
\end{proposition}

\begin{proof}
(i) Soit $\overline Z$ l'adhérence de $Z$ dans $X$~; comme $y$ est adhérent à $y'$ et $g$ continue, on a $g(y)\in\overline Z$. Comme $Z=\overline Z\cap f^{-1}(y')$, il résulte du théorème de semi-continuité de Chevalley (13.1.3) appliqué à la restriction de $f$ au sous-préschéma fermé réduit de $X$ ayant $\overline Z$ pour espace sous-jacent, que l'on a
\[
  \dim_{g(y)}(\overline Z\cap f^{-1}(y))\geqslant\dim(Z)~;
\]
mais les composantes irréductibles de $\overline Z\cap f^{-1}(y)$ contenant $g(y)$ sont évidemment contenues dans $X_y^0$, d'où \emph{a fortiori} l'inégalité (15.6.1.1).

(ii) Le raisonnement de (i) montre en outre que si les deux membres de (15.6.1.1) sont égaux, on a nécessairement
\[
  \dim_{g(y)}(\overline Z\cap f^{-1}(y))=\dim_{g(y)}(X_y^0)~;
\]
si $X_y^0$ est irréductible, cela entraîne $\overline Z\cap f^{-1}(y)=X_y^0$, donc $X_y^0$ est contenu dans $\overline Z$, et \emph{a fortiori} dans $\overline{X_{y'}^0}$.
\end{proof}

\oldpage[IV-3]{237}
\begin{corollary}[15.6.2]
\label{IV.15.6.2-fr}
Avec les notations de (15.6.1), supposons en outre que pour tout $y\in Y$, $X_y^0$ soit irréductible. Alors la fonction $z\mapsto\dim(X_z^0)$ est semi-continue supérieurement dans $Y$. Si en outre cette fonction est constante dans un voisinage d'un point $y\in Y$, alors on a $X_y^0\subset\overline{X_{y'}^0}$ pour toute générisation $y'$ de $y$.
\end{corollary}

\begin{proof}
Soit $y$ un point de $Y$, et soit $U$ un voisinage ouvert affine de $g(y)$ dans $X$~; alors il y a un voisinage ouvert $V$ de $y$ dans $Y$ tel que $g(V)\subset U$, et pour tout $z\in V$, $U\cap X_z^0$ est dense dans $X_z^0$, donc $\dim(X_z^0)=\dim(U\cap X_z^0)$ (4.1.1.3). Remplaçant $X$ par $U$, on peut donc supposer que $f$ est un morphisme de type fini. On sait alors (9.7.10) que la fonction $z\mapsto\dim(X_z^0)$ est localement constructible, et les assertions du corollaire résultent donc de (15.6.1) et de ($\mathbf{0}_{\mathrm{III}}$, 9.3.4).
\end{proof}

\begin{proposition}[15.6.3]
\label{IV.15.6.3-fr}
Avec les notations de (15.6.1), supposons que pour tout $y\in Y$, $X_y^0$ soit géométriquement irréductible (4.5.2). Alors, pour tout $y\in Y$ tel que la fonction $z\mapsto\dim(X_z^0)$ soit constante dans un voisinage de $y$, $f$ est universellement ouvert aux points de $X_y^0$.
\end{proposition}

\begin{proof}
Appliquons le critère (14.3.7). Soit donc $h:Y'\to Y$ un morphisme, $Y'$ étant le spectre d'un anneau de valuation discrète, dont nous désignerons par $t$, $t'$ le point fermé et le point générique respectivement, et supposons que $h(t)=y$, de sorte que $y'=h(t')$ est une générisation de $y$. Posons $X'=X\times_Y Y'$, $f'=f_{(Y')}:X'\to Y'$, $g'=g_{(Y')}:Y'\to X'$. Avec les mêmes notations que dans (15.6.1), l'hypothèse faite sur les $X_y^0$ implique que
\[
  X_t^{\prime 0}=X_y^0\otimes_{\mathbf{k}(y)}\mathbf{k}(t)
  \quad\text{et}\quad
  X_{t'}^{\prime 0}=X_{y'}^0\otimes_{\mathbf{k}(y')}\mathbf{k}(t'),
\]
et que $X_t^{\prime 0}$ et $X_{t'}^{\prime 0}$ sont irréductibles (4.4.1)~; comme $\dim(X_y^0)=\dim(X_t^{\prime 0})$ et $\dim(X_{y'}^0)=\dim(X_{t'}^{\prime 0})$ (4.1.4), on voit qu'on est ramené à prouver la proposition pour $f'$ et $X_t^{\prime 0}$, autrement dit on peut se borner au cas où $Y$ est spectre d'un anneau de valuation discrète, $y$ son point fermé et $y'$ son point générique~; si $x'$ est le point générique de $X_{y'}^0$, il résulte alors de (15.6.2) que tout point de $X_y^0$ est spécialisation de $x'$, d'où la conclusion.
\end{proof}

\begin{proposition}[15.6.4]
\label{IV.15.6.4-fr}
Sous les hypothèses générales de (15.6.1), soit $X^0$ la réunion des $X_y^0$ lorsque $y$ parcourt $Y$. Si $y\in Y$ est tel que $X_y^0$ soit géométriquement réduit sur $\mathbf{k}(y)$ (4.6.2) et si $f$ est universellement ouvert en tout point de $X_y^0$, alors $X^0$ est un voisinage de $X_y^0$ dans $X$. En particulier, si $f$ est universellement ouvert et si, pour tout $y\in Y$, $X_y^0$ est géométriquement réduit sur $\mathbf{k}(y)$, $X^0$ est ouvert dans $X$.
\end{proposition}

\begin{proof}
Montrons d'abord qu'on peut se ramener au cas où $f$ est un morphisme \emph{de type fini}. Soit $x$ un point de $X_y^0$~; il résulte de (5.10.8.1) (avec $\mathfrak F=\{\varnothing\}$) qu'il y a une suite finie $(Z_i)_{0\leqslant i\leqslant n}$ de composantes irréductibles de $f^{-1}(y)$ telle que $g(y)\in Z_0$, $x\in Z_n$ et $Z_i\cap Z_{i+1}\ne\varnothing$ pour $0\leqslant i\leqslant n-1$~; les $Z_i$ étant irréductibles, il y a une suite finie $(U_i)$ ($1\leqslant i\leqslant n$) d'ouverts affines dans $f^{-1}(y)$ tels que $g(y)\in U_1$, $x\in U_n$, $U_i$ étant un voisinage affine d'un point $x_i$ de $Z_i\cap Z_{i-1}$ pour $1\leqslant i\leqslant n$. Il y a pour chaque $i$ un ouvert quasi-compact $W_i$ dans $X$ tel que $U_i=f^{-1}(y)\cap W_i$~; soit $W$ l'ouvert quasi-compact de $X$ réunion des $W_i$. Comme $g$ est continue, il y a un voisinage ouvert affine $V$ de $y$ dans $Y$ tel que $g(V)\subset W$~; si $X'=W\cap f^{-1}(V)$, la restriction de $f$ à $X'$ est de type fini~; en outre, si, pour tout $z\in V$, $X_z^{\prime 0}$ est la composante connexe de $X'\cap f^{-1}(z)$ contenant $g(z)$, on a $X_z^{\prime 0}\subset X_z^0$, et $x$ appartient à $X_y^{\prime 0}$. En effet, chacun des $U_i$ contient les deux ensembles irréductibles $U_i\cap Z_{i-1}$ et $U_i\cap Z_i$ qui se rencontrent, et les ensembles irréductibles $U_i\cap Z_{i-1}$ et $U_{i-1}\cap Z_{i-1}$

\oldpage[IV-3]{238}
se rencontrent pour $2\leqslant i\leqslant n$~; la réunion des $U_i\cap Z_{i-1}$ et des $U_i\cap Z_i$ pour $1\leqslant i\leqslant n$ est donc connexe et contient $x$ et $g(y)$. Comme $f|X'$ est universellement ouvert aux points de $X_y^{\prime 0}$, cela achève de prouver notre assertion, car si la proposition est prouvée lorsque $f$ est de type fini, la réunion $X^{\prime 0}$ des $X_z^{\prime 0}$ pour $z\in V$ sera un voisinage de $x$, et il en sera de même de $X^0$.

Supposons donc $f$ de type fini. Alors $X^0$ est localement constructible (9.7.12)~; il suffit par suite de montrer que $X^0$ contient toute générisation $x'$ de $x$ ($\mathbf{0}_{\mathrm{III}}$, 9.2.5). Posons $y'=f(x')$~; il existe un anneau de valuation discrète $A$ et un morphisme $u$ de $Y'=\operatorname{Spec}(A)$ dans $X$ tels que, si $z$ est le point fermé et $z'$ le point générique de $Y'$, on ait $u(z)=x$ et $u(z')=x'$ (\textbf{II}, 7.1.9)~; si $h=f\circ u$, on a donc $h(z)=y$ et $h(z')=y'$. Posons $X'=X\times_Y Y'$, $f'=f_{(Y')}:X'\to Y'$, qui est un morphisme de type fini universellement ouvert aux points de $f'^{-1}(z)$, et $g'=g_{(Y')}$, qui est une $Y'$-section de $X'$. Si $p_y:f'^{-1}(z)\to f^{-1}(y)$ (resp. $p_{y'}:f'^{-1}(z')\to f^{-1}(y')$) est la projection canonique, $p_y^{-1}(X_y^0)$ (resp. $p_{y'}^{-1}(X_{y'}^0)$) est la composante connexe de $f'^{-1}(z)$ (resp. $f'^{-1}(z')$) contenant $g'(z)$ (resp. $g'(z')$), car les $X_t^0$ sont géométriquement connexes (4.5.13) et il suffit d'appliquer (4.4.1). De plus, le morphisme $u$ correspond à une $Y'$-section $v$ de $X'$ telle que $p_y(v(y))=x$, $p_{y'}(v(y'))=x'$, de sorte que $v(y')$ est une générisation de $v(y)$, et il suffira donc de prouver que $v(y')$ appartient à $p_{y'}^{-1}(X_{y'}^0)$. Enfin, il est clair que $p_y^{-1}(X_y^0)$ est géométriquement réduit sur $\mathbf{k}(z)$.

En d'autres termes, on est ramené au cas où $Y$ est spectre d'un anneau de valuation discrète, $y$ son point fermé, $y'$ son point générique. Comme $f^{-1}(y)-X_y^0$ est alors fermé dans $X$, on peut, en remplaçant $X$ par l'ouvert $X-(f^{-1}(y)-X_y^0)$, supposer que $X_y^0=f^{-1}(y)$~; de même, on peut remplacer $X$ par l'ouvert, composante connexe de $X$, qui contient $g(Y)$, cette composante connexe contenant évidemment $X^0$~; autrement dit, on peut supposer $X$ connexe. La proposition sera alors établie si l'on montre que $X^0=X$, ou encore que $X_{y'}$ est connexe~; mais $f^{-1}(y')$ est géométriquement réduit sur $\mathbf{k}(y')$, et \emph{a fortiori} réduit~; comme en outre $f$ est ouvert aux points de $f^{-1}(y)$, on peut appliquer le lemme (15.5.6), d'où la conclusion.
\end{proof}

\begin{corollary}[15.6.5]
\label{IV.15.6.5-fr}
Soient $f:X\to Y$ un morphisme plat de présentation finie, $g$ une $Y$-section de $X$~; pour tout $y\in Y$, soit $X_y^0$ la composante connexe de $f^{-1}(y)$ contenant $g(y)$, et soit $X^0$ la réunion des $X_y^0$ quand $y$ parcourt $Y$. Alors, pour tout $y\in Y$ tel que $X_y^0$ soit géométriquement réduit sur $\mathbf{k}(y)$, $X^0$ est un voisinage de $X_y^0$ dans $X$. En particulier, si $f$ est réduit (6.8.1), $X^0$ est ouvert dans $X$.
\end{corollary}

\begin{proof}
La question étant locale sur $Y$, on peut supposer $Y=\operatorname{Spec}(A)$ affine. Il existe donc un sous-anneau noethérien $A_0$ et un morphisme plat de type fini $f_0:X_0\to Y_0=\operatorname{Spec}(A_0)$ tels que $X=X_0\times_{Y_0}Y$ et $f=(f_0)_{(Y)}$ (8.9.1 et 11.2.7)~; en outre, on peut supposer qu'il existe une $Y_0$-section $g_0$ de $X_0$ telle que $g=(g_0)_{(Y)}$ (8.8.2). Pour tout $y\in Y$, soit $y_0$ sa projection dans $Y_0$~; il résulte de (4.5.13) que $(X_0)_{y_0}^0$ est géométriquement connexe, donc, si $p:f^{-1}(y)\to f_0^{-1}(y_0)$ est la projection canonique, on a $X_y^0=p^{-1}((X_0)_{y_0}^0)$ (4.4.1)~; de plus, l'hypothèse que $X_y^0$ est géométriquement réduit sur $\mathbf{k}(y)$ entraîne que $(X_0)_{y_0}^0$ est géométriquement réduit sur $\mathbf{k}(y_0)$ (4.6.10). On est ainsi ramené au cas où on suppose

\oldpage[IV-3]{239}
en outre $Y$ noethérien~; comme $f$ est universellement ouvert (11.1.1), il suffit alors d'appliquer (15.6.4).
\end{proof}

\begin{proposition}[15.6.6]
\label{IV.15.6.6-fr}
Soit $f:X\to Y$ un morphisme tel que $Y$ soit localement noethérien et $f$ localement de type fini, ou que $f$ soit de présentation finie. Soit $g$ une $Y$-section de $X$~; pour tout $y\in Y$, soit $X_y^0$ la composante connexe de $f^{-1}(y)$ contenant $g(y)$~; enfin, soit $X^0$ la réunion des $X_y^0$ lorsque $y$ parcourt $Y$.

Supposons que, pour tout $y\in Y$, $X_y^0$ soit géométriquement irréductible. Alors :
\begin{enumerate}[label=\emph{(\roman*)}]
  \item La fonction $y\mapsto\dim(X_y^0)$ est semi-continue supérieurement dans $Y$.
  \item Soient $x_0\in X^0$, $y_0=f(x_0)$. Si la fonction $y\mapsto\dim(X_y^0)$ est constante dans un voisinage de $y_0$, il existe un voisinage ouvert $U$ de $y_0$ tel que $f$ soit universellement ouvert aux points de $X^0\cap f^{-1}(U)$. Si de plus $Y$ est réduit et $f^{-1}(y_0)$ géométriquement réduit sur $\mathbf{k}(y_0)$ au point $x_0$, $f$ est plat au point $x_0$.
  \item Inversement, supposons que $f$ soit universellement ouvert au point générique de $X_{y_0}^0$ et en outre que l'une des conditions suivantes soit vérifiée :
  \begin{enumerate}
    \item[$\alpha$)] La fibre $f^{-1}(y_0)$ est géométriquement réduite sur $\mathbf{k}(y_0)$ au point $g(y_0)$ et l'anneau $\mathcal O_{Y,y_0}$ est noethérien.
    \item[$\beta$)] Pour tout point générique $y'$ d'une composante irréductible de $Y$ contenant $y_0$, on a $\dim(f^{-1}(y'))=\dim(X_{y'}^0)$.
  \end{enumerate}
  Alors la fonction $y\mapsto\dim(X_y^0)$ est constante dans un voisinage de $y_0$.
  \item L'ensemble $W$ des points $x\in X^0$ tels que la fonction $y\mapsto\dim(X_y^0)$ soit constante dans un voisinage de $f(x)$ et que $X_{f(x)}^0$ soit géométriquement réduit sur $\mathbf{k}(f(x))$, est ouvert dans $X$.
\end{enumerate}
\end{proposition}

\begin{proof}
Les questions étant locales sur $Y$, on peut supposer que $Y=\operatorname{Spec}(A)$ est affine. Lorsque $f$ est de présentation finie, il existe donc un sous-anneau noethérien $A_1$ de $A$ et un morphisme de type fini $f_1:X_1\to Y_1=\operatorname{Spec}(A_1)$ tels que $X=X_1\times_{Y_1}Y$ et $f=(f_1)_{(Y)}$ (8.9.1)~; en outre, on peut supposer qu'il existe une $Y_1$-section $g_1$ de $X_1$ telle que $g=(g_1)_{(Y)}$ (8.8.2). Pour tout $y\in Y$, soit $y_1$ sa projection dans $Y_1$~; il résulte de (4.5.13) que $(X_1)_{y_1}^0$ est géométriquement connexe, donc, si $p:f^{-1}(y)\to f_1^{-1}(y_1)$ est la projection canonique, on a $X_y^0=p^{-1}((X_1)_{y_1}^0)$ (4.4.1). Notons d'autre part que l'on peut, en vertu de (9.7.11) et (9.7.12), appliquer (9.3.3) à la propriété d'être géométriquement irréductible~; on peut donc supposer $A_1$ choisi de telle sorte que $(X_1)_{y_1}^0$ soit géométriquement irréductible pour tout $y_1\in Y_1$. On a $\dim(X_y^0)=\dim((X_1)_{y_1}^0)$ (4.1.4), et en appliquant de nouveau (9.3.3) à la propriété d'avoir une dimension donnée (et utilisant (9.5.5) et (9.7.12)), on peut supposer (en restreignant au besoin $Y$ à un voisinage de $y_0$) que si $\dim(X_y^0)$ est constante dans $Y$, alors $\dim((X_1)_{y_1}^0)$ est constante dans $Y_1$. Si $Y$ est réduit, il en est de même de $Y_1$ puisque $A_1\subset A$~; d'autre part, si $f^{-1}(y_0)$ est géométriquement réduit au point $x_0$, $f_1^{-1}(y_{01})$ l'est au point $x_{01}$ ($x_{01}$ et $y_{01}$ étant les projections respectives dans $X_1$ et $Y_1$ de $x_0$ et $y_0$) (4.6.10).

Ces remarques montrent que pour démontrer (i) et (ii), on peut se borner au cas où $Y$ est noethérien et $f$ localement de type fini. Mais dans ce cas, l'assertion (i) résulte de (15.6.2) et la première assertion de (ii) de (15.6.3). D'autre part, si $Y$ est réduit et $f^{-1}(y_0)$ géométriquement réduit sur $\mathbf{k}(y_0)$ au point $x_0$ ($Y$ étant toujours supposé

\oldpage[IV-3]{240}
noethérien), comme $X_y^0$ est la seule composante irréductible de $f^{-1}(y)$ contenant $x$, on déduit de (15.6.3) et (15.2.3) que si $\dim(X_y^0)$ est constante dans un voisinage de $y_0$, $f$ est plat au point $x_0$.

Pour prouver (iii), remarquons d'abord que l'ensemble $X^0$ est localement constructible dans $X$ (9.7.12), donc, par (9.5.5), l'ensemble $E$ des $y\in Y$ tels que $\dim(X_y^0)=\dim(X_{y_0}^0)$ est localement constructible dans $Y$. Appliquons alors (1.10.1) à $E$~: il suffit de prouver (compte tenu de (i)) que $\dim(X_{y'}^0)=\dim(X_{y_0}^0)$ pour le point générique $y'$ d'une composante irréductible de $Y$ contenant $y_0$~; ceci permet déjà de supposer que $Y=\operatorname{Spec}(\mathcal O_{Y,y_0})$.

Si on est dans le cas $\alpha$), on se ramène aussitôt, en vertu de (\textbf{II}, 7.1.9), et utilisant (4.5.13) et (4.4.1) comme dans (15.6.4), au cas où $Y$ est spectre d'un anneau de valuation discrète, $y_0$ son point fermé et $y'$ son point générique. Mais alors les hypothèses entraînent en vertu de (15.2.2.1), que $f$ est \emph{plat} au point $g(y_0)$, donc aussi dans un voisinage $V$ de ce point dans $X$ (11.1.1). Pour démontrer notre assertion, on peut remplacer $X$ par $V$, car $V\cap X_{y'}^0$ n'est pas vide en vertu de (2.3.4), donc est un ouvert partout dense dans $X_{y'}^0$ et a par suite même dimension (4.1.1.3)~; d'ailleurs, comme $X_{y'}^0$ est aussi la composante connexe de $f^{-1}(y')$ contenant $g(y')$, on peut supposer $V$ pris tel que $V$ ne rencontre aucune autre composante irréductible de $f^{-1}(y_0)$ ni de $f^{-1}(y')$, autrement dit on peut supposer que $X^0=X$~; mais alors la conclusion résulte de (12.1.1, (i)), puisque par hypothèse $X_{y_0}^0$ est intègre.

Supposons maintenant qu'on soit dans le cas $\beta$). Appliquons cette fois (\textbf{II}, 7.1.4) de la même manière que (\textbf{II}, 7.1.9) dans le cas $\alpha$)~: on est alors ramené au cas où $Y$ est spectre d'un anneau de valuation (non nécessairement discrète), $y_0$ son point fermé et $y'$ son point générique. En vertu de (14.3.13), il existe une composante irréductible $Z$ de $X$ contenant $X_{y_0}^0$ et dominant $Y$, et telle que $\dim(Z\cap f^{-1}(y'))=\dim(X_{y_0}^0)$~; mais l'hypothèse $\beta$) et le fait que $\dim(Z\cap f^{-1}(y'))\leqslant\dim(f^{-1}(y'))$ montrent que $\dim(X_{y_0}^0)\leqslant\dim(X_{y'}^0)$, d'où $\dim(X_{y_0}^0)=\dim(X_{y'}^0)$ en vertu de (i).

Reste à prouver (iv). Notons d'abord que les ensembles envisagés ne changent pas lorsqu'on remplace $f$ par $f_{(Y_{\mathrm{red}})}:X\times_Y Y_{\mathrm{red}}\to Y_{\mathrm{red}}$, la projection $X\times_Y Y_{\mathrm{red}}\to X$ étant un homéomorphisme. Autrement dit, on peut supposer $Y$ réduit, et alors il résulte de (ii) que $f$ est \emph{plat} aux points de $W$. Considérons un point $x_0$ de $W$ et prouvons que $W$ est un voisinage de $x_0$~; procédant comme au début de la démonstration, et utilisant aussi (11.2.7), on peut de plus supposer $Y$ noethérien~; il résulte alors de (ii) et de (15.6.4) que $X^0$ est un voisinage de $x_0$ dans $X$, et de (12.1.1, (vii)) que $W$ est aussi un voisinage de $x_0$ dans $X$.
\end{proof}

\begin{corollary}[15.6.7]
\label{IV.15.6.7-fr}
Supposons vérifiées les conditions préliminaires de (15.6.6) sur $f$, et supposons en outre que pour tout $y\in Y$, $X_y^0$ soit géométriquement intègre sur $\mathbf{k}(y)$. Alors les conditions suivantes sont équivalentes :
\begin{enumerate}[label=\emph{\alph*)}]
  \item La fonction $y\mapsto\dim(X_y^0)$ est localement constante dans $Y$.
  \item Le morphisme $f_{(Y_{\mathrm{red}})}:X\times_Y Y_{\mathrm{red}}\to Y_{\mathrm{red}}$ déduit de $f$ par changement de base, est plat aux points de $X^0$.
\end{enumerate}

\oldpage[IV-3]{241}
En outre, ces conditions entraînent que $X^0$ est ouvert dans $X$. Enfin, lorsque $Y$ est localement noethérien, a) et b) sont aussi équivalentes à
\begin{enumerate}
  \item[\rm b$'$)] $f$ est universellement ouvert aux points de $X^0$.
\end{enumerate}
\end{corollary}

\begin{proof}
Le fait que a) entraîne b) résulte de (15.6.6, (ii)), ainsi que le fait que $X^0$ est alors ouvert dans $X$. Si b) est vérifiée, on peut se borner au cas où $Y$ est réduit et $f$ plat~; alors, on se ramène, comme au début de (15.6.6), et en utilisant en outre (11.2.7), au cas où $Y$ est noethérien, cas où la conclusion résulte de (15.6.6, (iii), cas $\alpha$)). L'équivalence de b) et b$'$) a déjà été prouvée lorsque $Y$ est localement noethérien (15.2.2.1), compte tenu de ce que le morphisme $X\times_Y Y_{\mathrm{red}}\to X$ est un homéomorphisme universel.
\end{proof}

\begin{proposition}[15.6.8]
\label{IV.15.6.8-fr}
Soient $f:X\to Y$ un morphisme séparé de présentation finie, $g$ une $Y$-section de $X$~; pour tout $y\in Y$, soit $X_y^0$ la composante connexe de $f^{-1}(y)$ qui contient $g(y)$, et soit $X^0$ la réunion des $X_y^0$ quand $y$ parcourt $Y$. Alors, si $y\in Y$ est tel que $X_y^0$ soit propre sur $\mathbf{k}(y)$, $X^0$ est propre sur $Y$ au point $y$ (i.e. (15.7.1), il existe un voisinage ouvert $U$ de $y$ dans $Y$ tel que $X^0\cap f^{-1}(U)$ soit fermé dans $f^{-1}(U)$ et propre sur $U$).
\end{proposition}

\begin{proof}
Procédant comme au début de (15.6.6) (où on utilise (8.10.5, (v)), et où on remplace (9.7.11) par (9.6.7)), on se ramène au cas où $Y$ est noethérien et $f$ de type fini. On sait alors (9.7.12) que $X^0$ est localement constructible dans $X$~; d'autre part, les $X_z^0$ sont géométriquement connexes (pour tout $z\in Y$) en vertu de (4.5.13)~; enfin, comme $g(Y)\subset X^0$, $X^0$ est universellement submersif sur $Y$ (15.7.8). Il suffit alors d'appliquer à $X^0$ le critère (15.7.9).
\end{proof}

\begin{corollary}[15.6.9]
\label{IV.15.6.9-fr}
Soit $f:X\to Y$ un morphisme séparé tel que $Y$ soit localement noethérien et $f$ localement de type fini, ou que $f$ soit de présentation finie. Soit $g$ une $Y$-section de $X$, et pour tout $y\in Y$, soit $X_y^0$ la composante connexe de $f^{-1}(y)$ contenant $g(y)$. Soit $y$ un point de $Y$ tel que $X_y^0$ soit propre sur $\mathbf{k}(y)$. Alors, pour toute générisation $y'$ de $y$, on a $\dim(X_y^0)\geqslant\dim(X_{y'}^0)$.
\end{corollary}

\begin{proof}
Supposons d'abord $f$ de présentation finie~; alors en remplaçant au besoin $Y$ par un voisinage de $y$, on peut supposer que $f|X^0$ est \emph{propre} (15.6.8), et il suffit d'appliquer (13.1.5) à cette restriction. Si $Y$ est localement noethérien et $f$ localement de type fini, $X$ est localement noethérien~; en outre l'hypothèse que $X_y^0$ soit propre sur $\mathbf{k}(y)$ entraîne que $X_y^0$ est noethérien~; il existe donc un ouvert noethérien $V\subset X$ contenant $X_y^0$~; par suite il y a un voisinage ouvert $W$ de $y$ dans $Y$ tel que $g(W)\subset V$. En outre, si $z$ est point maximal de $X_{y'}^0$, on peut supposer que $z\in V$. La restriction de $f$ à $V\cap f^{-1}(W)$ est alors un morphisme de type fini, et $g|W$ une $W$-section~; en appliquant à ce morphisme le résultat déjà obtenu, on voit que si $Z$ est la composante irréductible de $X_{y'}^0$ de point générique $z$, on a $\dim(Z)\leqslant\dim(X_y^0)$. Comme $z$ est un point maximal quelconque de $X_{y'}^0$, on a bien $\dim(X_{y'}^0)\leqslant\dim(X_y^0)$.
\end{proof}

\begin{remarks}[15.6.10]
\label{IV.15.6.10-fr}
(i) L'hypothèse supplémentaire sur l'existence d'une composante irréductible de $X_{y'}^0$ contenant $g(y')$ et de dimension égale à $\dim(X_{y'}^0)$, faite dans (15.6.1), n'est pas superflue, comme le montre l'exemple suivant. Soient $A$ un anneau de valuation discrète, $\pi$ une uniformisante de $A$, $K$ le corps des fractions de $A$, $k$ son corps résiduel. Prenons $Y=\operatorname{Spec}(A)$, et désignons par $y$ et $y'$ le point fermé et le point générique de $Y$. Considérons les deux $Y$-schémas suivants~: $X_1=\operatorname{Spec}(A[t])$, $X_2=\operatorname{Spec}(K[u,v])$, où $t,u,v$ sont des indéterminées (on notera que la projection de $X_2$ dans $Y$ est donc $\{y'\}$). Dans l'anneau $A[t]$, l'idéal principal $(\pi t-1)$ est \emph{maximal}, et correspond donc à un point fermé $x_1$ de $X_1$, qui se projette au point générique $y'$ de $Y$ et est tel que $\mathbf{k}(x_1)=K$.

\oldpage[IV-3]{242}
Considérons d'autre part le point fermé $x_2$ de $X_2$ correspondant à l'idéal maximal $\left(u-\frac{1}{\pi}\right)+(v)$ de $K[u,v]$~; on a aussi $\mathbf{k}(x_2)=K$. Comme on le verra au chap.~V, on peut donc \emph{recoller} $X_1$ et $X_2$ le long des deux sous-préschémas fermés $Z_1$, $Z_2$ de $X_1$, $X_2$ ayant respectivement $\{x_1\}$ et $\{x_2\}$ pour espaces sous-jacents, suivant l'unique $Y$-isomorphisme de $Z_1$ sur $Z_2$. Désignons par $X$ le $Y$-préschéma ainsi obtenu, par $x$ le point de $X$ image de $x_1$ et $x_2$. La «~section nulle~» $g$ de $X_1$ (\textbf{II}, 1.7.9) est encore alors une $Y$-section de $X$~; $f^{-1}(y)$ est égale à $(X_1)_y$, tandis que $f^{-1}(y')$ a deux composantes irréductibles, respectivement isomorphes à $(X_1)_{y'}^0$ et $(X_2)_{y'}^0$, ayant un unique point commun $x$, et de dimensions respectives 1 et 2. Voir cependant la prop.~(15.6.9).

(ii) Considérons le morphisme $f$ défini dans (12.2.3, (ii)), qui est propre et plat~; en outre, la restriction de $f$ à $X_1\cap f_0^{-1}(Y)$ est un isomorphisme, donc le morphisme réciproque $g:Y\to X$ est une $Y$-section de $X$. On a alors $\dim(X_{y_0})=1$ tandis que $\dim(X_y^0)=0$ pour $y\ne y_0$, bien que $X_y^0$ soit géométriquement irréductible pour tout $y\in Y$ (mais $X_{y_0}^0$ n'est pas réduit)~; on voit donc que dans (15.6.6, (iii)), on ne peut supprimer les hypothèses $\alpha$) et $\beta$). En outre, $X^0$ n'est pas un voisinage de $X_{y_0}^{0}$, ce qui prouve que dans (15.6.4), on ne peut se dispenser de l'hypothèse que $X_{y_0}^0$ est réduit.

(iii) Au chap.~VI, nous appliquerons les résultats précédents aux $Y$-préschémas \emph{en groupes} localement de type fini sur un préschéma localement noethérien $Y$. Si $G$ est un tel préschéma, il existe une $Y$-section canonique $g$, la «~section unité~», et l'on montre que $G_y^0$ est toujours géométriquement irréductible et que l'on a $\dim(G_y^0)=\dim(G_y)$, en posant $G_y=f^{-1}(y)$. Il résultera donc de (15.6.2) et (15.6.3) que la fonction $z\mapsto\dim(G_z)$ est semi-continue supérieurement, et que si elle est constante au voisinage d'un point $y$, $f$ est universellement ouvert aux points de $G_y^0$. Si de plus $G_y$ est géométriquement réduit sur $\mathbf{k}(y)$ en un de ses points, on montre que $G_y$ est \emph{lisse} (6.8.1) sur $\mathbf{k}(y)$ en tous ses points~; il résultera alors de (15.6.6) que si de plus $\mathcal O_y$ est réduit, alors $f$ est lisse en tous les points de $G_y^0$ (mais pas nécessairement en tous les points de $G_y$). Ces résultats s'appliqueront par exemple dans la théorie de schémas de Picard, où nous disposerons d'un théorème général assurant que, sous certaines conditions, la fonction $z\mapsto\dim(G_z)$ est localement constante. Remarquons que c'est en vue d'applications de cette nature que les énoncés tels que (15.6.1) sont donnés pour les morphismes \emph{localement} de type fini, et non seulement pour les morphismes de type fini.

On notera aussi que dans le cas d'un $Y$-préschéma en groupes $G$, l'hypothèse $\beta$) de (15.6.6) est toujours vérifiée.
\end{remarks}

% EGA IV-3 terminal continuation, printed pp. 242--255.
% Direct diplomatic French transcription from the NUMDAM authority.

\subsection{Appendice : Critères valuatifs de propreté locale}
\label{subsection:IV.15.7-fr}

Ce numéro donne des compléments au critère valuatif de propreté démontré dans (\textbf{II}, 7.3.10)~; il est indépendant du reste du \S~15.

\begin{definition}[15.7.1]
\label{IV.15.7.1-fr}
Soient $f:X\to Y$ un morphisme de préschémas, $y$ un point de $Y$. On dit que $f$ est propre au point $y$ s'il existe un voisinage ouvert $U$ de $y$ dans $Y$ tel que la restriction $f^{-1}(U)\to U$ de $f$ soit un morphisme propre. On dit qu'une partie $Z$ de $X$ est propre sur $Y$ au point $y$ s'il existe un voisinage ouvert $U$ de $y$ tel que $Z\cap f^{-1}(U)$ soit fermée dans $f^{-1}(U)$ et propre sur $U$ (\textbf{II}, 5.4.10) (ce qui revient à dire que pour tout sous-préschéma fermé de $X$ ayant $Z$ pour espace sous-jacent, la restriction $Z\to Y$ de $f$ à $Z$ est propre au point $y$).
\end{definition}

Soit $g:Y'\to Y$ un morphisme quelconque~; posons $X'=X\times_Y Y'$, $f'=f_{(Y')}$ et, si $p:X'\to X$ est la projection canonique, $Z'=p^{-1}(Z)$. Alors, si $Z$ est propre sur $Y$ au point $y$, $Z'$ est propre sur $Y'$ en tout point $y'$ au-dessus de $y$ (\textbf{II}, 5.4.2).

\begin{proposition}[15.7.2]
\label{IV.15.7.2-fr}
Soient $Y$ un préschéma intègre localement noethérien, $X$ un préschéma intègre, $f:X\to Y$ un morphisme séparé, dominant et de type fini, $y$ un point de $Y$. Les conditions suivantes sont équivalentes :
\begin{enumerate}[label=\emph{\alph*)}]
  \item $f$ est propre au point $y$.
  \item Si $Y_1=\operatorname{Spec}(\mathcal O_y)$, $X_1=X\times_Y Y_1$, le morphisme $f_1=f_{(Y_1)}:X_1\to Y_1$ est propre.
  \item Tout anneau de valuation discrète ayant $R(X)$ pour corps des fractions et dominant $\mathcal O$ domine un anneau local $\mathcal O_x$ de $X$ (auquel cas on a nécessairement $f(x)=y$).
\end{enumerate}
\end{proposition}

\oldpage[IV-3]{243}
\begin{proof}
Le fait que a) implique b) résulte de (\textbf{II}, 5.4.2), et l'implication b) $\Rightarrow$ c) résulte de (\textbf{II}, 7.3.10). Reste donc à montrer que c) entraîne a). La question étant locale sur $Y$, on peut supposer $Y$ affine, donc noethérien. En vertu du lemme de Chow (\textbf{II}, 5.6.1), il existe un préschéma intègre $X'$, un morphisme projectif $p:P\to Y$, une immersion ouverte dominante $j:X'\to P$, et un morphisme projectif et birationnel (donc surjectif) $g:X'\to X$ tels que le diagramme
\[
\xymatrix{
P\ar[d]_-{p} & X'\ar[l]_-{j}\ar[d]^-{g}\\
Y & X\ar[l]^-{f}
}
\]
soit commutatif. Comme $X'$ est intègre et $j$ dominant, $P$ est irréductible, et l'on peut, en remplaçant $P$ par $P_{\mathrm{red}}$, supposer $P$ intègre (\textbf{I}, 5.2.2 et \textbf{II}, 5.5.5, (vi)). Tout revient à prouver qu'il existe un voisinage ouvert $U$ de $y$ tel que $p^{-1}(U)\subset j(X')$, car alors la restriction de $j$ à
\[
j^{-1}\bigl(p^{-1}(U)\bigr)=g^{-1}\bigl(f^{-1}(U)\bigr)=V
\]
sera un isomorphisme sur $p^{-1}(U)$, donc la restriction $V\to U$ de $f\circ g$ sera propre, et comme la restriction $V\to f^{-1}(U)$ de $g$ est surjective, la restriction $f^{-1}(U)\to U$ de $f$ sera propre (\textbf{II}, 5.4.3). Comme $T=P-j(X')$ est fermé dans $P$, $p(T)$ est fermé dans $Y$, et il suffit de montrer que $y\notin p(T)$, ou encore que tout $z'\in P$ tel que $p(z')=y$ \emph{appartient à} $j(X')$. Or, le corps des fractions rationnelles $R(P)$ est égal à $R(X')$, donc à $R(X)$ par construction; comme on peut se borner au cas où $z'$ n'est pas le point générique de $P$, il existe un anneau de valuation discrète $A$ ayant $R(X)$ pour corps des fractions et dominant $\mathcal O_{z'}$ (\textbf{II}, 7.1.7); comme $p(z')=y$, $A$ domine aussi $\mathcal O_y$, donc par hypothèse il existe un $x\in X$ tel que $f(x)=y$ et que $A$ domine $\mathcal O_x$. Comme $g$ est propre, il existe $x'\in X'$ tel que $g(x')=x$ et que $A$ domine $\mathcal O_{x'}$ (\textbf{II}, 7.3.10); donc $z'$ et $j(x')$ sont deux points de $P$ dont les anneaux locaux $\mathcal O_z$ et $\mathcal O_{j(x')}=\mathcal O_{x'}$ sont apparentés (\textbf{I}, 8.1.4); comme $P$ est un schéma, cela entraîne $z'=j(x')$ (\textbf{I}, 8.2.2), ce qui achève la démonstration.
\end{proof}

\begin{remark}[15.7.3]
\label{IV.15.7.3-fr}
On peut dans (15.7.2) supprimer l'hypothèse que $Y$ est localement noethérien, en remplaçant dans la condition c) les anneaux de valuation discrète par des anneaux de valuation quelconque; la démonstration est alors inchangée, compte tenu de (\textbf{II}, 7.3.10).
\end{remark}

\begin{corollary}[15.7.4]
\label{IV.15.7.4-fr}
Soient $Y$ un préschéma noethérien, $f:X\to Y$ un morphisme séparé de type fini, $Z$ une partie de $X$ tel que les points maximaux de son adhérence $\overline Z$ appartiennent à $Z$ (ce qui est le cas lorsque $Z$ est fini, ou lorsque $Z$ est constructible ($\mathbf{0}_{\mathrm{III}}$, 9.2.2)). Soit $y$ un point de $Y$. Les conditions suivantes sont équivalentes :
\begin{enumerate}
  \item[\rm a)] L'ensemble $\overline Z$ est propre sur $Y$ au point $y$.
  \item[\rm b)] Si $Y_1=\operatorname{Spec}(\mathcal O_y)$, $X_1=X\times_Y Y_1$, et si $g_1:X_1\to X$ est la projection canonique et $Z_1=g_1^{-1}(Z)$, l'adhérence $\overline Z_1$ de $Z_1$ dans $X_1$ est propre sur $Y_1$.
  \item[\rm c)] Pour tout schéma $Y'=\operatorname{Spec}(A')$, où $A'$ est un anneau de valuation discrète, et tout morphisme $g:Y'\to Y$, tel que l'image $g(y')$ du point fermé $y'$ de $Y'$ soit $y$, on a la propriété suivante : posant $X'=X\times_Y Y'$, $f'=f_{(Y')}:X'\to Y'$, $g'=g_{(X)}:X'\to X$, $Z'=g'^{-1}(Z)$, et désignant par $\eta'$
\oldpage[IV-3]{244}
le point générique de $Y'$, alors, pour tout point $x'\in Z'\cap f'^{-1}(\eta')$ rationnel sur $\mathbf k(\eta')$, il existe une $Y'$-section de $X'$ contenant $x'$.
\end{enumerate}
\end{corollary}

\begin{proof}
Le fait que a) entraîne b) résulte de ce que $\overline Z_1\subset g_1^{-1}(\overline Z)$ et de ce que $g_1^{-1}(\overline Z)$ est propre sur $Y_1$ (15.7.1). Le fait que b) entraîne c) résulte de (\textbf{II}, 7.3.3). Reste donc à voir que c) implique a). Il suffit évidemment de prouver que toute composante irréductible de $\overline Z$ est propre sur $Y$ au point $y$, et l'hypothèse permet donc de se borner au cas où $Z$ est réduit à un seul point $z$. On peut évidemment aussi, compte tenu de (\textbf{I}, 5.2.2) et (\textbf{II}, 5.4.6), supposer que $X$ et $Y$ sont réduits, puis (par définition d'une partie propre (\textbf{II}, 5.4.10)) supposer que $\overline Z=X=\overline{\{z\}}$ et (appliquant encore (\textbf{I}, 5.2.2)) que $f$ est dominant, donc $X$ et $Y$ intègres et noethériens; il faut alors prouver que $f$ est propre au point $y$ de $Y$. Montrons pour cela que $f$ vérifie la condition c) de (15.7.2). Soit $A'$ un anneau de valuation discrète ayant $R(X)=\mathbf k(z)$ pour corps des fractions et dominant $\mathcal O_y$; avec les notations de c), on a donc $g(y')=y$, et $g(\eta')$ est le point générique $\eta=f(z)$ de $Y$. Comme par définition $\mathbf k(\eta')=\mathbf k(z)$, il existe un $\mathbf k(\eta)$-morphisme $\operatorname{Spec}(\mathbf k(\eta'))\to\operatorname{Spec}(\mathbf k(z))$ transformant $\eta'$ en $z$, donc une $\mathbf k(\eta')$-section de $f'^{-1}(\eta')$ (\textbf{I}, 3.3.14); si $z'$ est l'image de $\eta'$ par cette section, $z'$ est donc rationnel sur $\mathbf k(\eta')$ et $g'(z')=z$. On conclut donc de l'hypothèse c) qu'il existe une $Y'$-section $h'$ de $X'$ telle que $h'(\eta')=z'$; soit $h=g'\circ h':Y'\to X$ le $Y$-morphisme correspondant, et soit $x=h(y')$; il est clair que $f(x)=y$ et que $A'=\mathcal O_{y'}$ domine $\mathcal O_x$, ce qui achève la démonstration.
\end{proof}

\begin{corollary}[15.7.5]
\label{IV.15.7.5-fr}
Soient $Y$ un préschéma localement noethérien, $f:X\to Y$ un morphisme séparé de type fini, $y$ un point de $Y$. Les conditions suivantes sont équivalentes :
\begin{enumerate}
  \item[\rm a)] $f$ est propre au point $y$.
  \item[\rm b)] Si $Y_1=\operatorname{Spec}(\mathcal O_y)$, $X_1=X\times_Y Y_1$, le morphisme $f_1=f_{(Y_1)}:X_1\to Y_1$ est propre.
  \item[\rm c)] Pour tout schéma $Y'=\operatorname{Spec}(A')$, où $A'$ est un anneau de valuation discrète, et tout morphisme $g:Y'\to Y$, tel que l'image $g(y')$ du point fermé $y'$ de $Y'$ soit $y$, on a la propriété suivante : posant $X'=X\times_Y Y'$, $f'=f_{(Y')}:X'\to Y'$ et désignant par $\eta'$ le point générique de $Y'$, alors, pour tout $x'\in f'^{-1}(\eta')$, rationnel sur $\mathbf k(\eta')$, il existe une $Y'$-section de $X'$ contenant $x'$.
\end{enumerate}
\end{corollary}

\begin{corollary}[15.7.6]
\label{IV.15.7.6-fr}
Soient $Y$ un préschéma localement noethérien, $f:X\to Y$ un morphisme de type fini. Les conditions suivantes sont équivalentes :
\begin{enumerate}
  \item[\rm a)] Il existe un morphisme propre $g:X\to Z$ et une immersion $h:Z\to Y$ tels que $f=h\circ g$.
  \item[\rm b)] Le sous-espace $f(X)$ est localement fermé dans $Y$, et si $Z'$ est le sous-préschéma de $Y$ image fermée de $X$ par $f$ (\textbf{I}, 9.5.3 et 9.5.1), et $Z''$ le préschéma induit par $Z'$ sur l'ouvert $f(X)$ de $Z'$, de sorte que $f$ se factorise de façon unique en $f=j\circ g$ (où $j:Z''\to Y$ est l'injection canonique), alors $g:X\to Z''$ est propre.
  \item[\rm c)] Pour tout $y\in f(X)$, $f$ est propre au point $y$.
  \item[\rm d)] Pour tout schéma $Y'=\operatorname{Spec}(A')$, où $A'$ est un anneau de valuation discrète, et tout morphisme $g:Y'\to Y$ tel que $g(Y')\subset f(X)$, toute $Y'$-section rationnelle (\textbf{I}, 7.1.2) de $X'=X\times_Y Y'$ se prolonge de façon unique en une $Y'$-section de $X'$.
\end{enumerate}
\end{corollary}

\begin{proof}
Il est clair que b) entraîne a). Pour voir que a) entraîne c), remarquons d'abord que a) entraîne que $h(Z)$ est localement fermé dans $Y$ et $g(X)$ fermé dans $Z$, donc $f(X)$ est localement fermé dans $Y$; pour tout $y\in f(X)$, il y a par suite un voisinage ouvert $U$
\oldpage[IV-3]{245}
de $y$ dans $Y$ tel que $f(X)\cap U$ soit fermé dans $U$; alors la restriction $h^{-1}(U)\to U$ de $h$ est une immersion fermée, et la restriction $f^{-1}(U)=g^{-1}(h^{-1}(U))\to h^{-1}(U)$ de $g$ est propre, donc la restriction $f^{-1}(U)\to U$ de $f$ est propre (\textbf{II}, 5.4.2). Pour voir que c) entraîne b), notons d'abord que, pour tout $y\in f(X)$, il existe, en vertu de c), un voisinage ouvert $U$ de $y$ dans $Y$ tel que $f(X)\cap U$ soit fermé dans $U$, donc $f(X)$ est localement fermé, et par suite ouvert dans son adhérence. Comme cette dernière est l'espace sous-jacent de $Z'$, et que $f$ se factorise en $f=j_0\circ g_0$, où $j_0:Z'\to Y$ est l'injection canonique et $g_0$ est un morphisme de $X$ dans $Z'$, le fait que $g_0(X)=Z''$ (qui est ouvert dans $Z'$) entraîne la factorisation de l'énoncé de b); le fait que $g$ est propre résulte alors du caractère local (sur $Z''$) de cette propriété, et de (\textbf{II}, 5.4.3). Il est clair que c) implique d) en vertu de (15.7.5), l'unicité du prolongement d'une $Y'$-section rationnelle résultant de l'hypothèse que $f$ est séparé (\textbf{I}, 7.2.2). Prouvons enfin que d) entraîne c); en premier lieu, il résulte de d), compte tenu de (\textbf{II}, 7.2.3 et 7.2.4), que $f$ est séparé; on peut alors appliquer le critère (15.7.5, c)), qui montre que $f$ est propre en tout point de $f(X)$.
\end{proof}

\begin{remarks}[15.7.7]
\label{IV.15.7.7-fr}
(i) Dans (15.7.4, c)), (15.7.5, c)) et (15.7.6, d)), on peut se restreindre au cas où l'anneau de valuation discrète $A'$ est \emph{complet} et admet un corps résiduel \emph{algébriquement clos}. En effet, dans la démonstration de (15.7.4), si l'on considère un anneau de valuation discrète $A''$ complet dominant $A'$ et ayant un corps résiduel algébriquement clos ($\mathbf{0}_{\mathrm{III}}$, 10.3.1), l'hypothèse c) est alors vérifiée pour $Y''=\operatorname{Spec}(A'')$, $X''=X\times_Y Y''$ et $f''=f_{(Y'')}$; mais $X''=X'\times_{Y'}Y''$, et il résulte alors de la remarque (\textbf{II}, 7.3.9, (i)) que $f'$ est propre; par suite $Y'$, $X'$ et $f'$ vérifient aussi l'hypothèse c) (\textbf{II}, 7.3.8).

(ii) Compte tenu de (\textbf{II}, 7.3.8), la condition c) de (15.7.5) peut se remplacer par la condition que $f'$ est \emph{propre}.
\end{remarks}

\begin{env}[15.7.8]
\label{IV.15.7.8-fr}
On dit qu'un morphisme de préschémas $f:X\to Y$ est \emph{submersif} s'il est surjectif et si la topologie de $Y$ est égale au quotient de la topologie de $X$ par la relation d'équivalence définie par $f$. On dit que $f$ est \emph{universellement submersif} si pour tout changement de base $Y'\to Y$, le morphisme $f'=f_{(Y')}:X\times_Y Y'\to Y'$ est submersif. Tout morphisme surjectif \emph{universellement ouvert}, ou \emph{universellement fermé}, ou \emph{fidèlement plat et quasi-compact} est universellement submersif (2.3.12). Étant donné un morphisme $f:X\to Y$, on dit qu'un sous-ensemble $Z$ de $X$ est \emph{submersif sur $f(Z)$} si la topologie du sous-espace $f(Z)$ de $Y$ est quotient de la topologie du sous-espace $Z$ de $X$ par la relation d'équivalence définie par $f|Z$. On dit que $Z$ est \emph{universellement submersif sur $f(Z)$} si, pour tout changement de base $g:Y'\to Y$, l'image réciproque $Z'=p^{-1}(Z)$ de $Z$ par la projection $p:X\times_Y Y'\to Y'$ est un ensemble submersif sur $f'(Z')=g^{-1}(f(Z))$. On notera que si le morphisme $f$ admet une $Y$-section $h:Y\to X$, tout ensemble $Z$ contenant $h(Y)$ est \emph{universellement submersif sur $Y$}.
\end{env}

\begin{proposition}[15.7.9]
\label{IV.15.7.9-fr}
Soient $Y$ un préschéma localement noethérien, $y$ un point de $Y$, $f:X\to Y$ un morphisme séparé de type fini. Soit $Z$ une partie localement constructible de $X$ ayant les propriétés suivantes :
\begin{enumerate}
  \item[\rm (i)] Pour toute générisation $y'$ de $y$, $Z_{y'}=Z\cap f^{-1}(y')$ est une composante connexe de $f^{-1}(y')$ et est géométriquement connexe sur $\mathbf k(y')$ (4.5.2).
\oldpage[IV-3]{246}
  \item[\rm (ii)] $Z_y$ est propre sur $\operatorname{Spec}(\mathbf k(y))$.
  \item[\rm (iii)] $Z$ est universellement submersif sur $f(Z)$ (15.7.8).
\end{enumerate}
Alors $Z$ est propre sur $Y$ au point $y$ (autrement dit (15.7.1) il existe un voisinage ouvert $U$ de $y$ tel que $Z\cap f^{-1}(U)$ soit fermé dans $f^{-1}(U)$ et propre sur $U$).
\end{proposition}

\begin{proof}
Utilisant la méthode de (8.1.2, a)) et (8.10.5, (xii)), on est ramené à prouver que lorsque $Y=\operatorname{Spec}(A)$ est spectre d'un anneau \emph{local} noethérien, les hypothèses (i), (ii) et (iii) entraînent que $Z$ est \emph{propre} sur $Y$.

Notons d'abord que si $A'$ est un anneau local noethérien, $\varphi:A\to A'$ un homomorphisme local, $Y'=\operatorname{Spec}(A')$ et $g:Y'\to Y$ le morphisme correspondant à $\varphi$, l'image réciproque $Z'$ de $Z$ par la projection canonique $p:X\times_Y Y'\to X$ a vis-à-vis de $Y'$ les propriétés (i), (ii), (iii) de l'énoncé, compte tenu de (1.8.2); utilisant (2.7.1, (vii)), il revient donc au même de démontrer que $Z'$ est propre sur $Y'$, pourvu que $A'$ soit un $A$-module \emph{fidèlement plat}. Nous allons utiliser cette remarque en prenant $A'=\widehat A$ ($\mathbf{0}_{\mathrm{I}}$, 7.3.5), autrement dit nous pouvons nous borner au cas où $A$ est \emph{complet}. Comme $Z_y$ est une composante connexe de $f^{-1}(y)$, propre sur $\mathbf k(y)$, il résulte de (\textbf{III}, 5.5.2) que $X$ est \emph{somme} de deux sous-préschémas $X_0$, $X_1$ induits sur des parties ouvertes et fermées de $X$, telles que $X_0$ soit propre sur $Y$ et que $X_0\cap f^{-1}(y)=Z_y$. Posons $Z_0=X_0\cap Z$, $Z_1=X_1\cap Z$, de sorte que ces ensembles forment une partition de $Z$ en deux parties ouvertes et fermées dans $Z$. Comme les $Z_{y'}$ sont connexes, on a nécessairement $Z_{y'}\subset Z_0$ ou $Z_{y'}\subset Z_1$, donc $f(Z_0)\cap f(Z_1)=\varnothing$. Mais comme $Z_0$ et $Z_1$ sont saturés pour la relation d'équivalence définie par $f|Z$, il résulte de (iii) que $f(Z_0)$ et $f(Z_1)$ sont ouverts dans $f(Z)$. Mais $f(Z_0)$ contient $y$, et tout voisinage de $y$ dans $Y$ est égal à $Y$, donc $f(Z_0)=f(Z)$, et par suite $f(Z_1)=\varnothing$, donc $Z_1=\varnothing$ et $Z\subset X_0$.

On peut donc désormais supposer de plus que $f$ est \emph{propre}, et il reste à prouver que $Z$ est fermé dans $X$. Autrement dit, il suffira de prouver qu'une partie \emph{constructible} $Z$ d'un préschéma $X$ \emph{propre} sur $Y$ qui satisfait aux \emph{deux seules conditions} (i) et (iii), est \emph{fermée} dans $X$. En vertu de ($\mathbf{0}_{\mathrm{III}}$, 9.2.5), il suffit donc de prouver que pour tout $x'\in Z$ et toute spécialisation $x$ de $x'$ dans $X$, on a $x\in Z$. Soient $A_1$ un anneau de valuation discrète \emph{complet}, $u$ un morphisme de $Y_1=\operatorname{Spec}(A_1)$ dans $X$ tel que, si $y_1$ et $y_1'$ sont le point fermé et le point générique de $Y_1$, on ait $u(y_1)=x$, $u(y_1')=x'$ (\textbf{II}, 7.1.7); posons $g=f\circ u:Y_1\to Y$ et $X_1=X\times_Y Y_1$; il existe une $Y_1$-section $u_1:Y_1\to X_1$ de $f_1=f_{(Y_1)}$ telle que $u=p\circ u_1$, où $p:X_1\to X$ est la projection canonique. Si $x_1=u_1(y_1)$, $x_1'=u_1(y_1')$, on a donc $p(x_1)=x$ et $p(x_1')=x'$, et $x_1$ est une spécialisation de $x_1'$. En outre, on a $x_1'\in Z_1=p^{-1}(Z)$; comme nous avons remarqué que les conditions (i) et (iii) sont stables pour \emph{tout} changement de base, ainsi que la propriété d'être constructible, on voit qu'on est ramené à la situation suivante : $Y$ est le spectre d'un anneau de valuation discrète complet, $X$ est propre sur $Y$, $f(x)=y$ est le point fermé de $Y$, $y'=f(x')$ son point générique, il existe une $Y$-section $h$ de $X$ telle que $h(y)=x$, $h(y')=x'$, et enfin on a $x'\in Z$; il faut prouver que $x\in Z$. Or, $Z_y$ est une composante connexe de $f^{-1}(y)$, \emph{propre} sur $\operatorname{Spec}(\mathbf k(y))$ puisque $X$ est propre sur $Y$. Appliquant de nouveau (\textbf{III}, 5.5.2), on voit comme plus haut qu'il y a une partie ouverte et fermée $X'$ de $X$ telle que $X'\cap f^{-1}(y)=Z_y$; comme $Z_y$ est connexe,
\oldpage[IV-3]{247}
on peut supposer $X'$ connexe (en remplaçant au besoin $X'$ par celle de ses composantes connexes contenant $Z_y$). Mais alors le même raisonnement qu'au début prouve qu'on a nécessairement $Z\subset X'$; comme $h(Y)$ est connexe et contient $x'$, il est nécessairement contenu dans $X'$, donc $x=h(y)\in X'\cap f^{-1}(y)=Z_y$. C.Q.F.D.
\end{proof}

\begin{corollary}[15.7.10]
\label{IV.15.7.10-fr}
Soient $Y$ un préschéma localement noethérien, $f:X\to Y$ un morphisme séparé de type fini, universellement submersif (15.7.8) et tel que toutes les fibres $X_y=f^{-1}(y)$ soient géométriquement connexes. Alors, si $y\in Y$ est tel que $X_y$ soit propre sur $\mathbf k(y)$, $f$ est propre au point $y$.
\end{corollary}

\begin{proof}
Compte tenu de (15.7.5), on est ramené au cas où $Y=\operatorname{Spec}(\mathcal O_y)$ puisque les hypothèses sont stables par changement de base. Mais dans ce cas il suffit d'appliquer (15.7.9) à $Z=X$.
\end{proof}

\begin{corollary}[15.7.11]
\label{IV.15.7.11-fr}
Soit $f:X\to Y$ un morphisme séparé de présentation finie; supposons qu'il existe un morphisme surjectif de présentation finie $g:Y'\to Y$, qui soit en outre propre ou plat, et tel que si l'on pose $X'=X\times_Y Y'$, il existe une $Y'$-section de $X'$ (ou encore un $Y$-morphisme $Y'\to X$ (\textbf{I}, 3.3.14)). Supposons enfin que les fibres $X_y=f^{-1}(y)$ soient géométriquement connexes. Alors l'ensemble $U$ des $y\in Y$ tels que $X_y$ soit propre sur $\mathbf k(y)$ est ouvert dans $Y$, et la restriction $f^{-1}(U)\to U$ de $f$ est propre.
\end{corollary}

\begin{proof}
La question étant locale sur $Y$, on peut supposer $Y=\operatorname{Spec}(A)$ affine. Appliquant (8.9.1) à $f$ et à $g$, on voit qu'il existe un sous-anneau noethérien $A_0$ de $A$ et deux morphismes de type fini $f_0:X_0\to Y_0=\operatorname{Spec}(A_0)$, $g_0:Y_0'\to Y_0$ tels que $X=X_0\times_{Y_0}Y$, $Y'=Y_0'\times_{Y_0}Y$, $f=(f_0)_{(Y)}$ et $g=(g_0)_{(Y)}$; en outre, utilisant (8.10.5, (v), (vi) et (xii)) et (11.2.7), on peut supposer $f_0$ séparé, $g_0$ surjectif et propre (resp. plat) si $g$ est propre (resp. plat); utilisant (8.8.2), on peut de plus supposer qu'il existe une $Y_0'$-section de $X_0$. D'autre part, utilisant (9.7.7), on peut appliquer (9.3.3) à la propriété d'être géométriquement connexe, et on peut donc supposer $A_0$ choisi de sorte que les $f_0^{-1}(y_0)$ soient géométriquement connexes pour tout $y_0\in Y_0$. Enfin, utilisant (2.7.1, (vii)), il revient au même de dire que $f^{-1}(y)$ est propre sur $\mathbf k(y)$ ou que $f_0^{-1}(y_0)$ est propre sur $\mathbf k(y_0)$, où $y_0$ est la projection de $y$ dans $Y_0$. Ces remarques montrent qu'on est ramené à démontrer le corollaire lorsque $Y$ est \emph{noethérien}. Notons maintenant le

\begin{lemma}[15.7.11.1]
\label{IV.15.7.11.1-fr}
Soient $f:X\to Y$ un morphisme, $g:Y'\to Y$ un morphisme surjectif, qui est en outre propre ou plat et quasi-compact. Posons $X'=X\times_Y Y'$, $f'=f_{(Y')}:X'\to Y'$. Alors, si $f'$ est submersif (resp. universellement submersif), il en est de même de $f$.
\end{lemma}

\begin{proof}
On peut se borner au cas où $f'$ est submersif. En effet, supposons le lemme prouvé dans ce cas, et montrons que si $f'$ est universellement submersif, $f$ l'est aussi. Soit donc $Y_1\to Y$ un morphisme quelconque, et posons $Y_1'=Y'\times_Y Y_1$; si $X_1'=X'\times_{Y'}Y_1'$ et $f_1'=(f')_{(Y_1')}:X_1'\to Y_1'$, $f_1'$ est submersif par hypothèse; de plus le morphisme $g_1:Y_1'\to Y_1$ est surjectif, et propre (resp. plat et quasi-compact) si $g$ l'est. Donc, si l'on pose $X_1=X\times_Y Y_1$, $f_1=f_{(Y_1)}$, il résulte de l'hypothèse et du fait que $X_1'=X_1\times_{Y_1}Y_1'$, que $f_1$ est submersif, donc $f$ universellement submersif. Supposons donc seulement que $f'$ soit submersif. Il est clair tout d'abord que $f$ est surjectif. Soit $E$ une partie de $Y$ telle que $f^{-1}(E)$ soit fermé dans $X$; alors, si $p:X'\to X$ est la projection canonique,
\oldpage[IV-3]{248}
$p^{-1}(f^{-1}(E))=f'^{-1}(g^{-1}(E))$ est fermé dans $X'$, donc, puisque $f'$ est submersif, $g^{-1}(E)$ est fermé dans $Y'$; mais comme $g$ est aussi universellement submersif (15.7.8), $E$ est fermé dans $Y$, ce qui prouve le lemme.
\end{proof}

Cela étant, comme il existe par hypothèse une $Y'$-section de $X'$, $f'$ est universellement submersif (15.7.8), donc il en est de même de $f$ d'après le lemme (15.7.11.1). Il suffit alors d'appliquer (15.7.10), en remarquant que l'ensemble des $y\in Y$ tels que $f$ soit propre en $y$ est par définition ouvert dans $Y$.
\end{proof}

\begin{center}
\emph{(A suivre.)}
\end{center}

\clearpage
\oldpage[IV-3]{249}
\phantomsection
\addcontentsline{toc}{section}{Bibliographie (suite)}
\section*{BIBLIOGRAPHIE (\emph{suite})}

\begingroup
\small
\setlength{\parindent}{0pt}
\setlength{\parskip}{3pt}
[39] J.-P.~\textsc{Serre}, \emph{Corps locaux}, Paris (Hermann), 1962.

[40] J.-P.~\textsc{Serre}, \emph{Groupes proalgébriques}, \emph{Publ. Inst. Hautes Études Sci.}, n\textsuperscript{o} 7 (1960).

\endgroup

\clearpage
\oldpage[IV-3]{250}
\phantomsection
\addcontentsline{toc}{section}{Index des notations}
\section*{INDEX DES NOTATIONS}

\begingroup
\small
\setlength{\parindent}{0pt}
\setlength{\parskip}{2pt}

$\mathfrak C(X)$, $\mathfrak D_c(X)$, $\mathfrak F_c(X)$,
$\mathfrak{LF}_c(X)$ ($X$ préschéma) : 8.3.9.

$\mathfrak Q(\mathcal F)$ ($\mathcal F$ $\mathcal O_X$-Module) : 8.5.10.

$\mathfrak{Spr}(Y)$, $\mathfrak{Spro}(Y)$, $\mathfrak{Sprf}(Y)$
($Y$ préschéma) : 8.6.1.

$\mathbf{Pro}(\mathcal C)$ ($\mathcal C$ catégorie) : 8.13.3.

$X_s$, $\mathcal F_s$, $u_s$, $g_s$, $Z_s$, $h_s$
($X$ $S$-préschéma, $s\in S$, $\mathcal F$ $\mathcal O_X$-Module,
$u$ homomorphisme de $\mathcal O_X$-Modules, $g$ section de
$\mathcal O_X$-Module, $Z$ partie de $X$, $h:X\to Y$ $S$-morphisme) : 9.4.1.

$\mathfrak{Qc}(X)$, $\mathfrak{Lqc}(X)$, $\mathfrak O(X)$,
$\mathfrak F(X)$, $\mathfrak{Lf}(X)$, $\mathfrak{Lc}(X)$
($X$ espace topologique) : 10.1.1.

$\operatorname{prof}^{*}_{x}(\mathcal F)$,
$\operatorname{prof}^{*}_{Z}(\mathcal F)$,
$\operatorname{prof}^{*}(\mathcal F)$
($\mathcal F$ $\mathcal O_X$-Module) : 10.8.1.

$S(X)$ ($X$ préschéma de Jacobson) : 10.9.1.

$S(f)$ ($f$ morphisme de préschémas de Jacobson) : 10.9.2.

$\operatorname{Spm}(A)$, $D^m(h)$ ($A$ anneau de Jacobson, $h\in A$) : 10.9.3.

$\operatorname{Tor.dim}_B(M)$ ($M$ $B$-module) : 12.3.2.

\endgroup

\clearpage
\oldpage[IV-3]{251}
\phantomsection
\addcontentsline{toc}{section}{Index terminologique}
\section*{INDEX TERMINOLOGIQUE}

\begingroup
\footnotesize
\setlength{\parindent}{0pt}
\setlength{\parskip}{1.1pt}
\newcommand{\EGAIVthreeTerm}[2]{#1 : #2.\par}

\EGAIVthreeTerm{Critère de platitude par fibres}{11.3.10}
\EGAIVthreeTerm{Espace algébrique sur $k$, espace $k$-algébrique, espace
préalgébrique sur $k$, espace $k$-préalgébrique}{10.10.2}
\EGAIVthreeTerm{Espace de Jacobson}{10.3.1}
\EGAIVthreeTerm{Espace $k$-préalgébrique réduit}{10.10.4}
\EGAIVthreeTerm{Famille de morphismes schématiquement dominante}{11.10.2}
\EGAIVthreeTerm{Famille de morphismes universellement schématiquement dominante}{11.10.7}
\EGAIVthreeTerm{Famille de sous-préschémas schématiquement dense}{11.10.2}
\EGAIVthreeTerm{Famille séparante d'homomorphismes de faisceaux de modules}{11.9.1}
\EGAIVthreeTerm{Famille séparante d'homomorphismes de modules}{11.9.4}
\EGAIVthreeTerm{Famille universellement séparante d'homomorphismes de faisceaux de modules}{11.9.14}
\EGAIVthreeTerm{Géométriquement isomorphes ($k$-préschémas)}{9.1.4}
\EGAIVthreeTerm{Homomorphisme de Modules universellement injectif}{11.9.14 et 11.9.18}
\EGAIVthreeTerm{Lemme de Chow pour les morphismes de présentation finie}{8.10.5.1}
\EGAIVthreeTerm{« Main theorem » de Zariski pour les morphismes de présentation finie}{8.12.6}
\EGAIVthreeTerm{Morphisme équidimensionnel en un point, morphisme équidimensionnel}{13.2.2 et 13.3.2}
\EGAIVthreeTerm{Morphisme propre en un point}{15.7.1}
\EGAIVthreeTerm{Morphisme pseudo-fini}{8.12.3}
\EGAIVthreeTerm{Morphisme submersif, morphisme universellement submersif}{15.7.8}
\EGAIVthreeTerm{Morphisme universellement ouvert en un point}{14.3.3}
\EGAIVthreeTerm{Ouvert ultraaffine}{10.9.5}
\EGAIVthreeTerm{Partie finie saturée d'un $k$-préschéma algébrique}{9.8.9}
\EGAIVthreeTerm{Partie propre en un point (d'un $S$-préschéma)}{15.7.1}
\EGAIVthreeTerm{Partie quasi-constructible, partie localement quasi-constructible d'un espace topologique}{10.1.1}
\EGAIVthreeTerm{Partie très dense d'un espace topologique}{10.1.3}
\EGAIVthreeTerm{Polynôme géométriquement irréductible}{9.7.4}
\EGAIVthreeTerm{Préschéma de Jacobson}{10.4.1}
\EGAIVthreeTerm{Préschéma équidimensionnel sur un autre en un point, équidimensionnel sur un autre}{13.2.2 et 13.3.2}
\EGAIVthreeTerm{Préschéma essentiellement affine sur un autre}{8.13.4}
\EGAIVthreeTerm{Principe de l'extension finie}{9.1.1}
\EGAIVthreeTerm{Profondeur rectifiée}{10.8.1}
\EGAIVthreeTerm{Pro-objet, morphisme de pro-objets}{8.13.3}
\EGAIVthreeTerm{Pro-objet essentiellement affine}{8.13.4, 8.14.1}
\EGAIVthreeTerm{Propriété constructible, propriété ind-constructible}{9.2.1 et 9.2.2, (i) et (ii)}
\EGAIVthreeTerm{Quasi-homéomorphisme d'espaces topologiques}{10.2.2}
\EGAIVthreeTerm{Quasi-isomorphisme d'espaces annelés}{10.2.8}
\EGAIVthreeTerm{Saturée d'une partie finie d'un $k$-préschéma algébrique}{9.8.9}
\EGAIVthreeTerm{Sous-ensemble d'un $S$-préschéma submersif, universellement submersif sur son image}{15.7.8}
\EGAIVthreeTerm{Squelette primaire d'un Module, squelette virtuel}{9.8.9}
\EGAIVthreeTerm{Spectre maximal d'un anneau de Jacobson}{10.9.3}
\EGAIVthreeTerm{Type primaire d'un Module, type primaire géométrique d'un Module}{9.8.9}
\EGAIVthreeTerm{Ultrapréschéma, morphisme d'ultrapréschémas}{10.9.5}
\EGAIVthreeTerm{$R$-ultrapréschéma}{10.10.2}

\let\EGAIVthreeTerm\relax
\endgroup

\clearpage
% La page imprimée 252 est blanche.
\thispagestyle{empty}
\null

\clearpage
\oldpage[IV-3]{253}
\phantomsection
\addcontentsline{toc}{section}{Table des matières}
\section*{TABLE DES MATIÈRES}

\begingroup
\small
\setlength{\parindent}{0pt}
\setlength{\parskip}{1.5pt}
\newcommand{\EGAIVthreeContentsLine}[3]{\textbf{#1}\quad #2\unskip\nobreak\hspace{0.4em}\nobreak\dotfill\nobreak\mbox{#3}\par}

\EGAIVthreeContentsLine{CHAPITRE IV.}{\textbf{Étude locale des schémas et des morphismes de schémas} \emph{(suite)}}{5}
\EGAIVthreeContentsLine{\S~8.}{Limites projectives de préschémas}{5}
\EGAIVthreeContentsLine{8.1.}{Introduction}{5}
\EGAIVthreeContentsLine{8.2.}{Limites projectives de préschémas}{7}
\EGAIVthreeContentsLine{8.3.}{Parties constructibles dans une limite projective de préschémas}{12}
\EGAIVthreeContentsLine{8.4.}{Critères d'irréductibilité et de connexion pour les limites projectives de préschémas}{17}
\EGAIVthreeContentsLine{8.5.}{Modules de présentation finie sur une limite projective de préschémas}{19}
\EGAIVthreeContentsLine{8.6.}{Sous-préschémas de présentation finie d'une limite projective de préschémas}{25}
\EGAIVthreeContentsLine{8.7.}{Critères pour qu'une limite projective de préschémas soit un préschéma réduit (resp. intègre)}{27}
\EGAIVthreeContentsLine{8.8.}{Préschémas de présentation finie sur une limite projective de préschémas}{28}
\EGAIVthreeContentsLine{8.9.}{Premières applications à l'élimination des hypothèses noethériennes}{34}
\EGAIVthreeContentsLine{8.10.}{Propriétés de permanence des morphismes par passage à la limite projective}{36}
\EGAIVthreeContentsLine{8.11.}{Application aux morphismes quasi-finis}{41}
\EGAIVthreeContentsLine{8.12.}{Nouvelle démonstration et généralisation du « Main Theorem » de Zariski}{43}
\EGAIVthreeContentsLine{8.13.}{Traduction en termes de pro-objets}{49}
\EGAIVthreeContentsLine{8.14.}{Caractérisation d'un préschéma localement de présentation finie sur un autre, en termes du foncteur qu'il représente}{52}

\EGAIVthreeContentsLine{\S~9.}{Propriétés constructibles}{54}
\EGAIVthreeContentsLine{9.1.}{Le principe de l'extension finie}{54}
\EGAIVthreeContentsLine{9.2.}{Propriétés constructibles et ind-constructibles}{56}
\EGAIVthreeContentsLine{9.3.}{Propriétés constructibles de morphismes de préschémas algébriques}{60}

\clearpage
\oldpage[IV-3]{254}
\EGAIVthreeContentsLine{9.4.}{Constructibilité de certaines propriétés des Modules}{62}
\EGAIVthreeContentsLine{9.5.}{Constructibilité de propriétés topologiques}{67}
\EGAIVthreeContentsLine{9.6.}{Constructibilité de certaines propriétés des morphismes}{71}
\EGAIVthreeContentsLine{9.7.}{Constructibilité des propriétés de séparabilité, d'irréductibilité géométrique et de connexité géométrique}{76}
\EGAIVthreeContentsLine{9.8.}{Décomposition primaire au voisinage d'une fibre générique}{83}
\EGAIVthreeContentsLine{9.9.}{Constructibilité des propriétés locales des fibres}{88}

\EGAIVthreeContentsLine{\S~10.}{Préschémas de Jacobson}{95}
\EGAIVthreeContentsLine{10.1.}{Parties très denses d'un espace topologique}{95}
\EGAIVthreeContentsLine{10.2.}{Quasi-homéomorphismes}{97}
\EGAIVthreeContentsLine{10.3.}{Espaces de Jacobson}{101}
\EGAIVthreeContentsLine{10.4.}{Préschémas de Jacobson et anneaux de Jacobson}{101}
\EGAIVthreeContentsLine{10.5.}{Préschémas de Jacobson noethériens}{104}
\EGAIVthreeContentsLine{10.6.}{Dimension dans les préschémas de Jacobson}{107}
\EGAIVthreeContentsLine{10.7.}{Exemples et contre-exemples}{109}
\EGAIVthreeContentsLine{10.8.}{Profondeur rectifiée}{110}
\EGAIVthreeContentsLine{10.9.}{Spectres maximaux et ultrapréschémas}{112}
\EGAIVthreeContentsLine{10.10.}{Espaces algébriques de Serre}{114}

\EGAIVthreeContentsLine{\S~11.}{Propriétés topologiques des morphismes plats de présentation finie. Critères de platitude}{116}
\EGAIVthreeContentsLine{11.1.}{Ensembles de platitude (cas noethérien)}{117}
\EGAIVthreeContentsLine{11.2.}{Platitude d'une limite projective de préschémas}{119}
\EGAIVthreeContentsLine{11.3.}{Application à l'élimination d'hypothèses noethériennes}{132}
\EGAIVthreeContentsLine{11.4.}{Descente de la platitude par des morphismes quelconques : cas d'un préschéma de base artinien}{143}
\EGAIVthreeContentsLine{11.5.}{Descente de la platitude par des morphismes quelconques : cas général}{150}
\EGAIVthreeContentsLine{11.6.}{Descente de la platitude par des morphismes quelconques : cas d'un préschéma de base unibranche}{154}
\EGAIVthreeContentsLine{11.7.}{Contre-exemples}{157}
\EGAIVthreeContentsLine{11.8.}{Un critère valuatif de platitude}{159}
\EGAIVthreeContentsLine{11.9.}{Familles séparantes et universellement séparantes d'homomorphismes de faisceaux de modules}{160}
\EGAIVthreeContentsLine{11.10.}{Familles schématiquement dominantes de morphismes et familles schématiquement denses de sous-préschémas}{170}

\EGAIVthreeContentsLine{\S~12.}{Étude des fibres des morphismes plats de présentation finie}{173}
\EGAIVthreeContentsLine{12.0.}{Introduction}{173}
\EGAIVthreeContentsLine{12.1.}{Propriétés locales des fibres d'un morphisme plat localement de présentation finie}{174}

\clearpage
\oldpage[IV-3]{255}
\EGAIVthreeContentsLine{12.2.}{Propriétés locales et globales des fibres d'un morphisme propre, plat et de présentation finie}{179}
\EGAIVthreeContentsLine{12.3.}{Propriétés cohomologiques locales des fibres d'un morphisme plat et localement de présentation finie}{183}

\EGAIVthreeContentsLine{\S~13.}{Morphismes équidimensionnels}{187}
\EGAIVthreeContentsLine{13.1.}{Le théorème de semi-continuité de Chevalley}{188}
\EGAIVthreeContentsLine{13.2.}{Morphismes équidimensionnels : cas des morphismes dominants de préschémas irréductibles}{190}
\EGAIVthreeContentsLine{13.3.}{Morphismes équidimensionnels : cas général}{194}

\EGAIVthreeContentsLine{\S~14.}{Morphismes universellement ouverts}{199}
\EGAIVthreeContentsLine{14.1.}{Morphismes ouverts}{200}
\EGAIVthreeContentsLine{14.2.}{Morphismes ouverts et formule des dimensions}{202}
\EGAIVthreeContentsLine{14.3.}{Morphismes universellement ouverts}{204}
\EGAIVthreeContentsLine{14.4.}{Le critère de Chevalley pour les morphismes universellement ouverts}{209}
\EGAIVthreeContentsLine{14.5.}{Morphismes universellement ouverts et quasi-sections}{216}

\EGAIVthreeContentsLine{\S~15.}{Étude des fibres d'un morphisme universellement ouvert}{223}
\EGAIVthreeContentsLine{15.1.}{Multiplicités des fibres d'un morphisme universellement ouvert}{223}
\EGAIVthreeContentsLine{15.2.}{Platitude des morphismes universellement ouverts à fibres géométriquement réduites}{226}
\EGAIVthreeContentsLine{15.3.}{Application : critères de réduction et d'irréductibilité}{228}
\EGAIVthreeContentsLine{15.4.}{Compléments sur les morphismes de Cohen-Macaulay}{229}
\EGAIVthreeContentsLine{15.5.}{Rang séparable des fibres d'un morphisme quasi-fini et universellement ouvert. Application aux composantes connexes géométriques des fibres d'un morphisme propre}{231}
\EGAIVthreeContentsLine{15.6.}{Composantes connexes des fibres le long d'une section}{236}
\EGAIVthreeContentsLine{15.7.}{Appendice : Critères valuatifs de propreté locale}{242}
\EGAIVthreeContentsLine{}{BIBLIOGRAPHIE \emph{(suite)}}{249}
\EGAIVthreeContentsLine{}{INDEX DES NOTATIONS}{250}
\EGAIVthreeContentsLine{}{INDEX TERMINOLOGIQUE}{251}

\vspace{2em}
\begin{center}
\emph{Manuscrit reçu le 15 novembre 1964.}
\end{center}

\let\EGAIVthreeContentsLine\relax
\endgroup

